\documentclass[11pt]{amsart}
\usepackage[a4paper, margin=1in]{geometry}
\usepackage{amsfonts}
\usepackage{amssymb}
\usepackage{amsmath}
\usepackage{amsthm}
\usepackage{mathtools}
\usepackage{hyperref}
\usepackage{cleveref}
\usepackage{stmaryrd}
\usepackage{tikz-cd}
\usepackage{enumitem}
\usepackage{multirow}
\usepackage{bbm}
\usepackage{comment}
\usepackage{float}
\usepackage{url}
%\usepackage{bm}

\title{Converse Theorems for Hilbert Modular Forms of Higher Level}
\author{Pengcheng Zhang}
\address{Max Planck Institute for Mathematics, Vivatsgasse 7, 53111 Bonn, Germany}
\email{pzhang@mpim-bonn.mpg.de}
\date{1 April, 2024}
\subjclass[2020]{11F41, 11F66}
\keywords{Hilbert modular forms, converse theorems, $L$-functions}

%\setlength\parindent{0pt}

\newtheorem{theorem}{Theorem}[section]

\newtheorem{defn}{Definition}

\newtheorem{lemma}[theorem]{Lemma}
\newtheorem{prop}[theorem]{Proposition}
\newtheorem{corol}[theorem]{Corollary}

\newtheorem{conj}{Conjecture}
\newtheorem{ques}{Question}

\theoremstyle{remark}
\newtheorem*{rmk}{Remark}

\newcommand\m[1]{\begin{pmatrix}#1\end{pmatrix}}
\newcommand\bm[1]{\begin{bmatrix}#1\end{bmatrix}}
\newcommand\sm[1]{\begin{psmallmatrix}#1\end{psmallmatrix}}
\newcommand\bsm[1]{\begin{bsmallmatrix}#1\end{bsmallmatrix}}
\newcommand\mbf[1]{\mathbf{#1}}
\newcommand\cl[1]{\mathcal{#1}}
\newcommand\fr[1]{\mathfrak{#1}}
\newcommand\ov[1]{\overline{#1}}
\newcommand\p[2]{\mathbb{P}^{#1}(#2)}
\newcommand\wt[1]{\widetilde{#1}}
\newcommand\wh[1]{\widehat{#1}}

\newcommand{\rat}{\mathbb{Q}}
\newcommand{\inte}{\mathbb{Z}}
\newcommand{\real}{\mathbb{R}}
\newcommand{\comp}{\mathbb{C}}
\newcommand{\ade}{\mathbb{A}}
\newcommand{\nat}{\mathbb{N}}
\newcommand{\alg}{\overline{\mathbb{Q}}}
\newcommand{\pos}{\mathbb{Z}^+}
\newcommand{\fp}{\mathbb{F}_p}
\newcommand{\fpn}{\mathbb{F}_{p^n}}
\newcommand{\qp}{\mathbb{Q}_p}
\newcommand{\oqp}{\overline{\mathbb{Q}_p}}
\newcommand{\zp}{\mathbb{Z}_p}
\newcommand{\cp}{\mathbb{C}_p}
\newcommand{\ql}{\mathbb{Q}_\ell}
\newcommand{\oql}{\overline{\mathbb{Q}_\ell}}
\newcommand{\eps}{\varepsilon}
\newcommand{\vphi}{\varphi}
\newcommand{\clh}{\mathcal{H}}
\newcommand{\cld}{\mathcal{D}}
\newcommand{\clo}{\mathcal{O}}
\newcommand{\clu}{\mathcal{U}}
\newcommand{\clm}{\mathcal{M}}
\newcommand{\clf}{\mathcal{F}}
\newcommand{\dprime}{\prime\prime}
\newcommand{\sbullet}{\scriptscriptstyle\bullet}
\newcommand{\blank}{\hspace{0.5cm}}
\newcommand{\bslash}{\backslash}
\newcommand{\simarrow}{\xrightarrow{\;\sim\;}}
\newcommand{\eq}{\;=\;}
\newcommand{\deq}{\;:=\;}
\newcommand{\bone}{\mathbf{1}}
\def\thin{{\hskip 1pt}}
\newcommand{\hash}{\texttt{\#}}

\newcommand{\suml}{\sum\limits}
\newcommand{\prodl}{\prod\limits}
\newcommand{\unionl}{\bigcup\limits}
\newcommand{\disjunl}{\bigsqcup\limits}
\newcommand{\interl}{\bigcap\limits}
\newcommand{\liml}{\lim\limits}
\newcommand{\osuml}{\bigoplus\limits}

\newcommand{\id}{\mathrm{id}}
\newcommand{\chr}{\mathrm{char}}
\newcommand{\Hom}{\mathrm{Hom}}
\newcommand{\End}{\mathrm{End}}
\newcommand{\Aut}{\mathrm{Aut}}
\newcommand{\M}{\mathrm{M}}
\newcommand{\N}{\mathrm{N}}
\newcommand{\GL}{\mathrm{GL}}
\newcommand{\PGL}{\mathrm{PGL}}
\newcommand{\PSL}{\mathrm{PSL}}
\newcommand{\SL}{\mathrm{SL}}
\newcommand{\SO}{\mathrm{SO}}
\newcommand{\SU}{\mathrm{SU}}
\newcommand{\Mp}{\mathrm{Mp}}
\newcommand{\tr}{\mathrm{tr}}
\newcommand{\Gal}{\mathrm{Gal}}
\newcommand{\Stab}{\mathrm{Stab}}
\newcommand{\Div}{\mathrm{Div}}
\newcommand{\dv}{\mathrm{div}}
\newcommand{\ord}{\mathrm{ord}}
\newcommand{\Ker}{\mathrm{Ker}}
\newcommand{\tors}{\mathrm{tors}}
\newcommand{\Res}{\mathrm{Res}}
\newcommand{\dr}{\mathrm{dR}}
\newcommand{\Corr}{\mathrm{Corr}}
\newcommand{\Fr}{\mathrm{Fr}}
\newcommand{\Gmod}{G\textit{-}\mathrm{mod}}
\newcommand{\Ab}{\mathrm{Ab}}
\newcommand{\ab}{\mathrm{ab}}
\newcommand{\Ext}{\mathrm{Ext}}
\newcommand{\para}{\mathrm{par}}
\newcommand{\Art}{\mathrm{Art}}
\newcommand{\ur}{\mathrm{ur}}
\newcommand{\tame}{\mathrm{tame}}
\newcommand{\lcm}{\mathrm{lcm}}
\newcommand{\et}{\mathrm{\acute{e}t}}
\newcommand{\val}{\mathrm{val}}
\newcommand{\Cl}{\mathrm{Cl}}
\newcommand{\sgn}{\mathrm{sgn}}
\renewcommand{\Re}{\operatorname{Re}}
\renewcommand{\Im}{\operatorname{Im}}

\newcommand{\SG}{\mathrm{S}\Gamma}

\begin{document}
\maketitle

\begin{abstract}
We prove two results on converse theorems for Hilbert modular forms of higher level over totally real fields of degree $r>1$. The first result recovers a Hilbert modular form (of some level) from an $L$-series satisfying functional equations twisted by all the Hecke characters that are unramified at all finite places. The second result assumes both the above functional equations and an Euler product, and recovers a Hilbert modular form of the expected level predicted by the shape of the functional equations.
\end{abstract}

\setcounter{tocdepth}{1}
\tableofcontents

\section{Introduction and the main result}

\subsection{Background and introduction}

The study of converse theorems, i.e., recovering automorphic forms from $L$-series, dates back to Hecke \cite{hecke}. He proved an equivalence between $L$-series satisfying a suitable functional equation and functions invariant under the element $\sm{0&-1\\1&0}$, which is then sufficient to recover elliptic modularity of level $1$. Hecke's work was generalized by Weil \cite{weil}, who recovered elliptic modularity of general level $N$ by assuming functional equations twisted by all the Dirichlet characters. Later in \cite{conreyfarmer}, Conrey and Farmer turned to another direction of generalization by imposing an Euler product to the $L$-series and recovered the modularity for a few small levels. 

Converse theorems for Hilbert modular forms was first investigated by Doi--Naganuma \cite{doinaganuma} and Naganuma \cite{naganuma} when they discovered a lifting from elliptic modular forms to Hilbert modular forms, nowadays known as the Doi--Naganuma lifting. In order to justify that the image of the lifting is a Hilbert modular form, they proved a converse theorem for Hilbert modular forms of level $1$ and applied it to the lifting, but subject to the base field being norm-Euclidean. In fact, subject to the class number $1$ condition, they proved an equivalence between $L$-series satisfying functional equations twisted by all the unramified Hecke characters and functions invariant under the element $\sm{0&-1\\1&0}$, but they could not justify that $\sm{0&-1\\1&0}$ and the translation matrices always generate the full Hilbert modular group. It was later proved by Vasersteion \cite{vaserstein} that indeed this generating result always holds.

In the case of general number fields, it is well-known that Jacquet and Langlands \cite{jacquetlanglands} proved a converse theorem for $\GL_2$ in the style of Weil, i.e., assuming functional equations twisted by all the Hecke characters. 

The goal of this article is to recover Hilbert modularity of higher level from $L$-series with suitable properties. Specifically, we will be focused on the following two properties of an $L$-series:
\begin{enumerate}[label=(\arabic*)]
    \item having suitable functional equations twisted by all the unramified Hecke characters;
    \item having a suitable (partial) Euler product.
\end{enumerate}
We will prove that one can already recover a Hilbert modular form (of some level) by assuming all the `unramified' functional equations, and together with the Euler product, one can recover a Hilbert modular form of the expected level predicted by the shape of the functional equations.

Roughly speaking, we will associate the $L$-series $L(s)$ with an $h$-tuple $\mbf{f}=(f_1,\cdots,f_h)$ of holomorphic functions on $\clh^r$ whose $L$-function equals $L(s)$, and prove that $\mbf{f}$ is a Hilbert modular form by showing that each component $f_\lambda$ is a classical Hilbert modular form. Here $h$ is the narrow class number of the base field and $r>1$ is the degree of the base field. The proof can be divided into four steps.
\begin{enumerate}[leftmargin=1.5cm, label=Step \arabic*.]
    \item Modify the proof in \cite[Section~3]{doinaganuma} to deduce from the `unramified' functional equations of $L(s)$ that $\mbf{f}$  is an eigenfunction of the `Fricke involution' (\Cref{fqnsection}).
    \item With $\mbf{f}$ now being an eigenfunction of the `Fricke involution', use the main theorem of Vaserstein in \cite{vaserstein} to deduce that each $f_\lambda$ is already invariant under a congruence subgroup of a larger level, which proves the first part of the theorem (\Cref{generating-sect,gamma1-inv}).
    \item Use the particular shape of the Euler factor of $L(s)$ at certain primes to deduce that each $f_\lambda$ is an eigenfunction of the `Hecke operator at infinity' at that prime (\Cref{eulerprodsection}).
    \item With $f_\lambda$ now being a `Hecke eigenfunction' at those certain primes, use the idea from \cite[Theorem~2.2]{conjhecke} to extend the congruence subgroup that $f_\lambda$ is invariant under to the `correct' one predicted by the shape of the functional equations, which proves the second part of the theorem (\Cref{gamma0-inv-sect}).
\end{enumerate}

We would like to point out that there is a small but easily fixable gap in the proof of \cite[Theorem~2.2]{conjhecke}. We will also provide our own fix at the end of \Cref{gamma0-inv-sect}. For the coherence of the discussion, we postpone certain lemmas on prime ideals to the appendix.

\subsection{Setup and the main result}
\label{statement}
Let $F$ be a totally real field of degree $r>1$ and let $\clo_F$ be the ring of integers. Let $\sigma_1,\cdots,\sigma_r:F\rightarrow\real$ be the real embeddings of $F$. For $\alpha\in F$, write $\alpha^{(j)}=\sigma_j(\alpha)$, and $\alpha$ is called \textit{totally positive} if $\alpha^{(j)}>0$ for all $1\leq j\leq r$. For any set $E\subseteq F$, let $E_+$ denote the set of totally positive elements in~$E$. In particular, we write $\clo_+=(\clo_F)_+$ and $\clo^\times_+=(\clo_F^\times)_+$ for convenience. Let $\cl{I}_F$ denote the set of nonzero integral ideals of $\clo_F$ and let $\fr{d}$ denote the different of $F$. Let $\N$ denote both the norm map $\N_{F/\rat}$ on the elements of $F$ and the norm map on the (fractional) ideals of $F$. Let $\Cl^+(F)$ denote the narrow class group of $F$ and let $h=|\Cl^+(F)|$ denote the narrow class number of $F$. For any nonzero (fractional) ideal $\fr{a}$ of~$F$, let $[\fr{a}]$ denote the equivalence class of $\fr{a}$ in $\Cl^+(F)$. We will view each equivalence class in $\Cl^+(F)$ as a subset of $\cl{I}_F$. Let $\{\fr{t}_1,\cdots,\fr{t}_h\}\subseteq\cl{I}_F$ be a fixed set of representatives of $\Cl^+(F)$. 

We will give a brief discussion of Hecke characters and the classical formulation of Hilbert modular forms (of non-parallel weight) in \Cref{prelims}. Throughout the article, unramified Hecke characters mean Hecke characters that are unramified at every finite places. For convenience, we will use the notations in \Cref{prelims} without defining them here.

Let $\{A(\fr{a})\}_{\fr{a}\in \cl{I}_F}$ be a sequence of complex numbers indexed by nonzero integral ideals of $\clo_F$ such that $A(\clo_F)=1$ and $A(\fr{a})=O(\N(\fr{a})^c)$ for some constant $c>0$. Define an $L$-series
\begin{align*}
    L(s)\eq L(s,A)\deq\sum_{\fr{a}\in\cl{I}_F}A(\fr{a})\cdot \N(\fr{a})^{-s}.
\end{align*}
We will usually omit the $A$ in $L(s,A)$ and in the later notations when there is no confusion, as we think of $\{A(\fr{a})\}_{\fr{a}\in \cl{I}_F}$ as a fixed sequence.

For a nonzero integral ideal $\fr{m}$ of $\clo_F$, define
\begin{align*}
    K_0(\fr{m})&\deq\bigg\{\m{a&b\\c&d}\in\GL_2(\wh{\clo_F})\;\bigg|\;c\in\fr{m}\bigg\} \\
    K_1(\fr{m})&\deq\bigg\{\m{a&b\\c&d}\in\GL_2(\wh{\clo_F})\;\bigg|\;c\in\fr{m},\;d-1\in\fr{m}\bigg\},
\end{align*}
where
\begin{align*}
    \wh{\clo_F}\deq\varprojlim_{\fr{t}}\clo_F/\fr{t}\eq\prod_{\fr{p}}\clo_{F,\fr{p}}.
\end{align*}

Now, fix a nonzero integral ideal $\fr{n}$ of $\clo_F$. Fix an $r$-tuple $\mbf{k}=(k_1,\cdots,k_r)\in(2\inte^+)^r$ of positive even integers. Write $k_0=\max_{1\leq j\leq r}k_j$, $k_j^\prime=k_0-k_j$, and $\mbf{k}^\prime=(k_1^\prime,\cdots,k_r^\prime)$. For each Hecke character~$\psi$ of $F$, define a twisted $L$-series
\begin{align*}
L(s,A,\psi)\deq\sum_{\fr{a}\in\cl{I}_F}\psi(\fr{a})\cdot A(\fr{a})\cdot \N(\fr{a})^{-s}
\end{align*}
and define the completed $L$-series as
\begin{align*}
    \wh{L}(s,A,\psi)\deq(2\pi)^{-rs}\; \N(\fr{n}\fr{d}^2)^{s/2}\thin\prod_{j=1}^r\Gamma\bigg(s-\frac{k_j^\prime}{2}+\nu_j(\psi)\bigg)\thin L(s,A,\psi),
\end{align*}
where we will define $\nu_j(\psi)$ in \Cref{hecke-char-sect}.

\begin{theorem}
\label{mainthm}
Let $F$ be a totally real field of degree $r>1$. Let $\fr{n}$ be a nonzero integral ideal of~$\clo_F$, $\mbf{k}\in(2\inte^+)^r$, and $\epsilon=\pm1$. Let $\{\fr{t}_1,\cdots,\fr{t}_h\}$ be any set of representatives of $\Cl^+(F)$ and let $\fr{m}=\fr{t}_1\fr{t}_2\cdots\fr{t}_h$. Suppose that for all the Hecke characters $\psi$ of $F$ that are unramified at all finite places, all the functions $\wh{L}(s,A,\psi)$ have holomorphic continuation to the whole complex plane which are bounded in vertical strips and satisfy functional equations
\begin{align*}
    \wh{L}(s,A,\psi)\eq \epsilon \thin i^{-|\mbf{k}|}\;\psi(\fr{n}\fr{d}^2)\;\wh{L}(k_0-s,A,\ov{\psi}),
\end{align*}
where $|\mbf{k}|=\sum_{j=1}^rk_j$. Then, $L(s)$ is the $L$-function of a Hilbert modular form of weight $\mbf{k}$ and level $K_1(\fr{m}\fr{n})$.

Suppose that, moreover, the $L$-series $L(s)$ has an Euler product of the form
\begin{align*}
    L(s)\eq\prod_{\fr{p}\text{ prime}}L_\fr{p}(s),
\end{align*}
where for $\fr{p}\nmid\fr{n}\fr{d}^2$,
\begin{align*}
    L_\fr{p}(s)\eq\big(1-A(\fr{p})\N(\fr{p})^s+\N(\fr{p})^{k_0-1-2s}\big)^{-1}.
\end{align*}
Then, $L(s)$ is the $L$-function of a Hilbert modular form of weight $\mbf{k}$ and level $K_0(\fr{n})$.
\end{theorem}
\begin{rmk}
We make several remarks on the theorem.
\begin{enumerate}[label=(\roman*)]
    \item If $F$ has narrow class number $1$, then $\fr{m}=\clo_F$ in the first part of the theorem, so all the `unramified' functional equations already recover a Hilbert modular form of the expected level. The two parts of the theorem can thus be viewed as a trade-off between having narrow class number~$1$ and having an Euler product, i.e., if one does not want to assume one of them, then one has to assume the other in order to obtain the expected level.
    \item In some sense, the case when $F$ has narrow class number $1$ can be viewed as a special case of the main theorem in \cite{p-s}. Piatetski-Shapiro's theorem implies a converse theorem for $\GL_2$ assuming all the `unramified' functional equations when the base field has class number $1$ and at least one real place, but his theorem is written in the language of automorphic representations. One should also turn to \cite[(iii),~p.671]{bookerkrish-strengthen} and \cite[Lemma~2.2]{bookerkrish-refine} for a more detailed discussion.
    \item For the second part of the theorem, we in fact only need a partial Euler product at primes $\fr{p}$ that are of degree $1$, with $[\fr{p}]=[\clo_F]$ in $\Cl^+(F)$, and $\fr{p}\nmid\fr{n}\fr{d}^2$.
    \item This theorem should be viewed as a generalization of the converse theorem proved in Doi--Naganuma \cite{doinaganuma} and Naganuma \cite{naganuma}. Indeed, with our theorem, one may be able to follow Doi--Naganuma's ideas to prove the Hilbert modularity of the image of the Doi--Naganuma lifting in more general cases.
    \item In the case of $\GL_{n(\geq3)}$ over number fields, Cogdell and Piatetski-Shapiro \cite[Theorem~3]{cogdellpsi} first proved a similar converse theorem in the sense of only assuming functional equations twisted by all the unramified automorphic representations of lower ranks, but subject to the class number $1$ condition. The class number $1$ restriction was later removed by Booker and Krishnamurthy \cite[Theorem~1.1]{bookerkrish-gln}. In this sense, our theorem may be viewed as a partial extension of this result from $\GL_{n(\geq3)}$ to $\GL_2$.
\end{enumerate}
\end{rmk}

\subsection{Acknowledgements}
The author would like to thank his supervisor, Don Zagier, for all the advice he has brought up regarding both mathematics and writing. The author would also like to thank Andrew Booker for useful remarks and correspondences. This work was carried out at Max Planck Institute for Mathematics.


\section{Preliminaries and notations}
\label{prelims}
In this section, we will review the theory of Hecke characters and the classical formulation of Hilbert modular forms. We will also write down a candidate which we will later prove is a Hilbert modular form whose $L$-function equals $L(s)$. Our main reference for Hecke characters is Bump's book \cite{bump}, and the one for Hilbert modular forms is Shimura's paper \cite{shimurahmf}, though in both cases there are some slight notational differences.

\subsection{Notations}

We first make a list of commonly used notations which we adopt for convenience. Let $\clh=\{z\in\comp\mid\Im(z)>0\}$ be the upper half plane. Define
\begin{align*}
    \GL_2^+(F)\deq\{\gamma\in\GL_2(F)\mid(\det\gamma)^{(j)}>0\text{ for all }1\leq j\leq r\}.
\end{align*}
Let $z=(z_1,\cdots,z_r)\in\clh^r$, $\kappa=(\kappa_1,\cdots,\kappa_r)\in\comp^r$, $\xi,\mu\in F$, and $\gamma\in\GL_2^+(F)$. We also denote $\gamma^{(j)}=\sm{a^{(j)}&b^{(j)}\\c^{(j)}&d^{(j)}}$ for $\gamma=\sm{a&b\\c&d}$.
We write
\begin{align*}
    \xi z+\mu&\eq(\xi^{(1)}z_1+\mu^{(1)},\cdots,\xi^{(r)}z_r+\mu^{(r)}) \\
    \gamma z&\eq(\gamma^{(1)}z_1,\cdots,\gamma^{(r)}z_r) \\
    z^\kappa&\eq\prod_{j=1}^rz_j^{\kappa_j},\hspace{0.5cm}\xi^\kappa\eq\prod_{j=1}^r(\xi^{(j)})^{\kappa_j} \\
    \tr(\xi z)&\eq\sum_{j=1}^r \xi^{(j)}z_j
\end{align*}
whenever they make sense.


\subsection{Hecke characters}
\label{hecke-char-sect}

Recall that a Hecke character of $F$ is a continuous group homomorphism $\psi:F^\times\bslash\ade_F^\times\rightarrow\comp^\times$, where $\ade_F$ is the adele ring of $F$. Let $\psi$ be a Hecke character of $F$. For each place $v$ of $F$, let $\psi_v$ denote the composition of $\psi$ and $F_v^\times\rightarrow F^\times\bslash\ade_F^\times$. We also define  $\psi_\infty:F_\infty=F\otimes_\rat\real\rightarrow\comp^\times$ as
\begin{align*}
    \psi_\infty\deq\prod_{j=1}^r\psi_{\sigma_j}
\end{align*}
where $\sigma_j$ are real places of $F$. It is known that at each real place $\sigma_j$ of $F$, there exist $\delta_j\in\{0,1\}$ and $\nu_j\in\comp$ such that
\begin{align*}
    \psi_{\sigma_j}(x^{(j)})\eq\sgn(x^{(j)})^{\delta_j}\thin|x^{(j)}|^{\nu_j}
\end{align*}
for all $x\in F$. Define $\nu_j(\psi)=\nu_j$ for each $1\leq j\leq r$.

Now, let $\psi$ be an unramified Hecke character of $F$, i.e., $\psi_v(\clo_{F_v}^\times)=1$ for all finite places $v$ of~$F$. For a fractional ideal $\fr{a}$ of $F$, we may define
\begin{align*}
    \psi(\fr{a})\deq\prod_{i=1}^n\psi_{\fr{p}_i}(\varpi_{\fr{p}_i})^{e_i},
\end{align*}
where $\fr{a}=\prod_{i=1}^n\fr{p}_i^{e_i}$ is the unique factorization of $\fr{a}$ into prime ideals and $\varpi_{\fr{p}_i}$ is a uniformizer of~$F_{\fr{p}_i}$. This is well-defined since $\psi$ is unramified at all finite places.

For the later purpose, we would like to construct a special family $(\psi_m)_{m\in\inte^{r-1}}$ of unramified Hecke characters of $F$. Fix a (multiplicative) basis $\{\eps_1,\cdots,\eps_{r-1}\}$ of~$\clo_+^\times$. For each \mbox{$m=(m_1,\cdots,m_{r-1})\in\inte^{r-1}$}, there exists (unique) \mbox{$\nu_m=(\nu_{m,1},\cdots,\nu_{m,r})\in (i\real)^r$} such that
\begin{equation}
\begin{aligned}
\label{nulinearrel}
    \sum_{j=1}^r\nu_{m,j}&\eq 0 \\
    \sum_{j=1}^r \nu_{m,j}\log (\eps_l^{(j)})&\eq 2\pi i m_l \hspace{0.5cm}\text{for }1\leq l\leq r-1.
\end{aligned}
\end{equation}
It is possible to make a choice of $(\delta_{m,j})_{j=1}^r\in\{0,1\}^r$ so that $\chi_m:F_\infty\rightarrow\comp$ given by
\begin{align*}
    \chi_m((x_j)_{j=1}^r)\deq\prod_{j=1}^r\sgn(x_j)^{\delta_{m,j}}\thin|x_j|^{\nu_{m,j}}
\end{align*}
defines a character on $\clo_F^\times\bslash F_\infty$ (see \cite[Section~1.7,~p.78]{bump}). Then, there exists a (not necessarily unique) Hecke character $\psi_m$ such that $(\psi_m)_{\infty}=\chi_m$ (see Exercise 5 in Section 6 of \cite[Chapter~VII]{neukirch}). In particular,
\begin{enumerate}[label=\rm{(\arabic*)}]
    \item $\nu_j(\psi_m)=\nu_{m,j}$ for all $1\leq j\leq r$;
    \item $\psi_m((\xi))=\prod_{j=1}^r(\xi^{(j)})^{-\nu_{m,j}}=\xi^{-\nu_m}$ for all totally positive $\xi\in F_+$.
\end{enumerate}
We will choose $(\psi_m)_{m\in\inte^{r-1}}$ so that $\psi_m\psi_{m^\prime}=\psi_{m+m^\prime}$. In particular, $\psi_{-m}=\psi_m^{-1}=\ov{\psi_m}$ and $\psi_{\mbf{0}}$ is the trivial Hecke character where $\mbf{0}=(0,\cdots,0)$. 

\subsection{Hilbert modular forms}

For $\gamma=\sm{a&b\\c&d}\in\GL_2^+(F)$ and \mbox{$\mbf{k}=(k_1,\cdots,k_r)\in(2\inte^+)^r$}, define the \textit{slash operator of weight $\mbf{k}$} on a holomorphic function $g:\clh^r\rightarrow\comp$ as
\begin{equation}
\begin{aligned}
\label{slashoperator}
    (g|_\mbf{k}\gamma)(z)&\deq\prod_{j=1}^r\big((\det\gamma^{(j)})^{k_j/2}\thin(c^{(j)}z_j+d^{(j)})^{-k_j}\big)\cdot g(\gamma^{(1)} z_1,\cdots,\gamma^{(r)}z_r) \\
    &\;\eq(\det\gamma)^{\mbf{k}/2}\cdot(cz+d)^{-\mbf{k}}\cdot g(\gamma z).
\end{aligned}
\end{equation}
Let $j:\GL_2^+(F)\rightarrow\PGL_2^+(F)$ be the projection map. A subgroup $\Gamma\subseteq\GL_2^+(F)$ is called an \textit{arithmetic} subgroup if $j(\Gamma\cap\GL_2^+(\clo_F))$ is a finite index subgroup of both $j(\Gamma)$ and $j(\GL_2^+(\clo_F))$. Let $\Gamma\subseteq\GL_2^+(F)$ be an arithmetic subgroup. A holomorphic function $g:\clh^r\rightarrow\comp$ is called a \textit{(classical) Hilbert modular form of weight $\mbf{k}$ and level $\Gamma$} if 
\begin{enumerate}[label=(\arabic*)]
    \item $g|_\mbf{k}\gamma=g$ for all $\gamma\in\Gamma$;
    \item for all $\gamma\in\GL_2^+(\clo_F)$, $g|_\mbf{k}\gamma$ is bounded as $\prod_{j=1}^r\Im(z_j)\rightarrow\infty$.
\end{enumerate}
It is well-known that when $r>1$, Koecher's principle implies that the second condition is not needed for the definition of Hilbert modular forms (e.g.~see \cite[Theorem~1.20]{123}).

We will be focused on adelic Hilbert modular forms on $K_0(\fr{n})$ and $K_1(\fr{n})$. We will not recall the definition of adelic Hilbert modular forms, but rather use the following equivalent viewpoint. Recall that $\{\fr{t}_1,\cdots,\fr{t}_h\}$ is a set of representatives of $\Cl^+(F)$. For two nonzero integral ideals $\fr{t}$ and $\fr{m}$ of $\clo_F$, define
\begin{align*}
    \Gamma_0^+(\fr{t},\fr{m})&\deq\bigg\{\m{a&b\\c&d}\in\GL_2^+(F)\;\bigg|\;a\in\clo_F,\;b\in\fr{t}^{-1},\;c\in\fr{t}\fr{m},\;d\in\clo_F,\;ad-bc\in\clo_+^\times\bigg\} \\
    \Gamma_1^+(\fr{t},\fr{m})&\deq\bigg\{\m{a&b\\c&d}\in\GL_2^+(F)\;\bigg|\;a\in\clo_F,\;b\in\fr{t}^{-1},\;c\in\fr{t}\fr{m},\;d-1\in\fr{m},\;ad-bc\in\clo_+^\times\bigg\},
\end{align*}
where $\fr{t}^{-1}:=\{a\in F\thin|\thin(a)\fr{t}\subseteq\clo_F\}$. (In the standard notations, the groups $\Gamma_i^+(\fr{t},\fr{m})$ are simply denoted as $\Gamma_i(\fr{t},\fr{m})$. Here we add the $+$ sign to maintain consistency with the notations in \Cref{generating-sect}.) Then, for $i=0,1$, an adelic Hilbert modular form of weight~$\mbf{k}$ and level $K_i(\fr{n})$ can be (equivalently) viewed as an $h$-tuple $\mbf{g}=\{g_1,\cdots,g_h\}$ where each $g_\lambda$ is a classical Hilbert modular form of weight $\mbf{k}$ and level $\Gamma_i^+(\fr{t}_\lambda,\fr{n})$ for $1\leq\lambda\leq h$.

Let \mbox{$\mbf{g}=(g_1,\cdots,g_h)$} be an adelic Hilbert modular form of weight $\mbf{k}$ and level $K_i(\fr{n})$. Note that $\Gamma_i^+(\fr{t}_\lambda,\fr{n})$ contains $T^\alpha=\sm{1&\alpha\\0&1}$ for all $\alpha\in\fr{t}_\lambda^{-1}$. The condition that $g_\lambda|_\mbf{k}T^\alpha=g_\lambda$ then implies that $g_\lambda(z+\alpha)=g_\lambda(z)$ for all $\alpha\in\fr{t}^{-1}$. By the theory of Fourier series, $g_\lambda$ has a Fourier series expansion of the form
\begin{align*}
    g_\lambda(z)\eq b_\lambda(0)+\sum_{\xi\in(\fr{t}_\lambda\fr{d}^{-1})_+}b_\lambda(\xi)\thin e(\tr(\xi z)),
\end{align*}
where $e(x)=e^{2\pi ix}$. Here we recall that the different $\fr{d}$ of $F$ is the integral ideal satisfying that
\begin{align*}
    \fr{d}^{-1}\eq\clo_F^\vee\deq\{x\in F\mid\tr(x\alpha)\in\inte\text{ for all }\alpha\in\clo_F\}.
\end{align*}
We will assume for simplicity that $b_\lambda(0)=0$ for all $1\leq\lambda\leq h$.

Note that $\Gamma_i^+(\fr{t}_\lambda,\fr{n})$ also contains $E_\eps=\sm{\eps&0\\0&1}$ for all $\eps\in\clo_+^\times$. By comparing the Fourier coefficients, the condition that $g_\lambda|_\mbf{k}E_\eps=g_\lambda$ then implies that the value $b_\lambda(\xi)\xi^{\mbf{k}^\prime/2}$ only depends on the coset $\xi\clo_+^\times$, where we write $k_0=\max_{1\leq j\leq r}k_j$, $k_j^\prime=k_0-k_j$, and $\mbf{k}^\prime=(k_1^\prime,\cdots,k_r^\prime)$ as before. The $L$-function of $g_\lambda$ can then be defined as
\begin{align*}
    L(s,g_\lambda)\deq\sum_{\xi\in(\fr{t}_\lambda\fr{d}^{-1})_+/\clo_+^\times}\N(\fr{t}_\lambda)^{k_0/2}\thin b_\lambda(\xi)\thin \xi^{\mbf{k}^\prime/2}\cdot\N(\xi\fr{t}_\lambda^{-1}\fr{d})^{-s}.
\end{align*}

Let $\cl{C}_\lambda:=[\fr{t}_\lambda^{-1}\fr{d}]$ denote an equivalence class in $\Cl^+(F)$. Note that every $\fr{a}\in\cl{C}_\lambda$ can be written as $\xi\fr{t}_\lambda^{-1}\fr{d}$ for some $\xi\in(\fr{t}_\lambda\fr{d}^{-1})_+$. Hence, the $L$-function of $g_\lambda$ can be rewritten as
\begin{align*}
    L(s,g_\lambda)\eq\sum_{\fr{a}\in\cl{C}_\lambda}B(\fr{a})\cdot\N(\fr{a})^{-s}.
\end{align*}
where $B(\fr{a})$ is defined via $B(\xi\fr{t}_\lambda^{-1}\fr{d}):=\N(\fr{t}_\lambda)^{k_0/2}b_\lambda(\xi)\xi^{\mbf{k}^\prime/2}$.

Now, define the $L$-function of $\mbf{g}=(g_1,\cdots,g_h)$ as
\begin{align*}
    L(s,\mbf{g})\deq\sum_{\fr{a}\in\cl{I}_F}B(\fr{a})\cdot\N(\fr{a})^{-s}\eq\sum_{\lambda=1}^hL(s,g_\lambda).
\end{align*}
For each Hecke character $\psi$, define the twisted $L$-function as
\begin{align*}
    L(s,\mbf{g},\psi)\deq\sum_{\fr{a}\in\cl{I}_F}\psi(\fr{a})\cdot B(\fr{a})\cdot\N(\fr{a})^{-s}
\end{align*}
and define the completed $L$-function as
\begin{align*}
    \wh{L}(s,\mbf{g},\psi)\deq(2\pi)^{-rs}\; \N(\fr{n}\fr{d}^2)^{s/2}\thin\prod_{j=1}^r\Gamma\bigg(s-\frac{k_j^\prime}{2}+\nu_j(\psi)\bigg)\thin L(s,\mbf{g},\psi).
\end{align*}

For each $1\leq\lambda\leq h$, let $\wt{\lambda}$ denote the unique integer such that $[\fr{t}_\lambda\fr{t}_{\wt{\lambda}}\fr{n}]=1$ in $\Cl^+(F)$. Let $q_\lambda$ be a totally positive generator of $\fr{t}_\lambda\fr{t}_{\wt{\lambda}}\fr{n}$ and define $W_{q_\lambda}=\sm{0&-1\\q_\lambda&0}$. Let $\epsilon=\pm1$. Suppose that $g_\lambda|_\mbf{k}W_{q_\lambda}=\epsilon g_{\wt{\lambda}}$ for all $1\leq\lambda\leq h$. Then, all the completed $L$-functions satisfy the functional equations
\begin{align*}
    \wh{L}(s,\mbf{g},\psi)\eq \epsilon\thin i^{-|\mbf{k}|}\;\psi(\fr{n}\fr{d}^2)\;\wh{L}(k_0-s,\mbf{g},\ov{\psi}).
\end{align*}
This is the motivation for the particular shape of functional equations considered in \Cref{mainthm}.

It is also possible to define the Hecke operators on $\mbf{g}$, but we will not go into details here as it is quite technical and complicated. One should turn to \cite[Section~2]{shimurahmf} for a detailed discussion. Suppose that $\mbf{g}$ is an eigenfunction of the Hecke operator at a prime $\fr{p}$ with $\fr{p}\nmid\fr{n}\fr{d}^2$. Then, the $L$-function of $\mbf{g}$ has a $\fr{p}$-partial Euler product of the form
\begin{align*}
    L(s,\mbf{g})\eq L^\fr{p}(s,\mbf{g})\;L_\fr{p}(s,\mbf{g})
\end{align*}
with
\begin{align*}
   L_\fr{p}(s,\mbf{g})\eq\big(1-B(\fr{p})\N(\fr{p})^s+\N(\fr{p})^{k_0-1-2s}\big)^{-1}\hspace{0.5cm}\text{and}\hspace{0.5cm}
    L^\fr{p}(s,\mbf{g})\eq\sum_{\fr{p}\nmid\fr{a}}B(\fr{a})\cdot\N(\fr{a})^{-s}.
\end{align*}
This is the motivation for the particular shape of the Euler product considered in \Cref{mainthm}.

\subsection{A candidate for the Hilbert modular form}

Now, we go back to the setting of \Cref{mainthm}. Define an $h$-tuple $\mbf{f}=(f_1,\cdots,f_h)$ as the following. For each $1\leq\lambda\leq h$, define a holomorphic function $f_\lambda:\clh^r\rightarrow\comp$ by
\begin{align*}
    f_\lambda(z)\deq\sum_{\xi\in(\fr{t}_\lambda\fr{d}^{-1})_+}a_\lambda(\xi)\thin e(\tr(\xi z)),
\end{align*}
where $a_\lambda(\xi)=\N(\fr{t}_\lambda)^{k_0/2}A(\xi\fr{t}_\lambda^{-1}\fr{d})\xi^{-\mbf{k}^\prime/2}$ for $\xi\in(\fr{t}_\lambda\fr{d}^{-1})_+$.

\begin{lemma}
\label{fhmf}
If $\mbf{f}=(f_1,\cdots,f_h)$ is an adelic Hilbert modular form of weight $\mbf{k}$ and level $K_i(\fr{m})$ for some nonzero ideal $\fr{m}$ of $\clo_F$, then \mbox{$L(s,\mbf{f})=L(s)$}, and in particular, $L(s)$ is an $L$-function of a Hilbert modular form of weight~$\mbf{k}$ and level~$K_i(\fr{m})$.
\end{lemma}
\begin{proof}
This follows easily from the previous review on Hilbert modular forms.
\end{proof}

By \Cref{fhmf}, we will be focused on proving that $\mbf{f}$ is an adelic Hilbert modular form from now on, i.e., each $f_\lambda$ is a classical Hilbert modular form of certain corresponding level. We will finish this section with a basic lemma on $f_\lambda$. Recall that $T^\alpha=\sm{1&\alpha\\0&1}$ for $\alpha\in F$ and that $E_\eps=\sm{\eps&0\\0&1}$ for $\eps\in\clo_+^\times$.

\begin{lemma}
\label{basic}
For each $1\leq\lambda\leq h$, $f_\lambda$ satisfies that $f_\lambda|_\mbf{k}T^\alpha=f_\lambda$ for all $\alpha\in\fr{t}_\lambda^{-1}$ and that $f_\lambda|_\mbf{k}E_\eps=f_\lambda$ for all  $\eps\in\clo_+^\times$.
\end{lemma}
\begin{proof}
It follows from the Fourier series expansion of $f_\lambda$ that $f_\lambda|_\mbf{k}T^\alpha=f_\lambda$ for all $\alpha\in\fr{t}_\lambda^{-1}$. For~$E_\eps$, we have
\begin{align*}
    (f_\lambda|_\mbf{k}{E_\eps})(z)&\eq\eps^{\mbf{k}/2}\thin f_\lambda(\eps z) \\
    &\eq\eps^{\mbf{k}/2}\sum_{\xi\in(\fr{t}_\lambda\fr{d}^{-1})_+}\N(\fr{t}_\lambda)^{k_0/2}\thin A(\xi\fr{t}_\lambda^{-1}\fr{d})\thin \xi^{-\mbf{k}^\prime/2}\thin e(\tr(\xi\eps z)) \\
    &\eq\eps^{\mbf{k}/2}\sum_{\xi\in(\fr{t}_\lambda\fr{d}^{-1})_+}\N(\fr{t}_\lambda)^{k_0/2}\thin A(\xi\eps^{-1}\fr{t}_\lambda^{-1}\fr{d})\thin (\xi\eps^{-1})^{-\mbf{k}^\prime/2}\thin e(\tr(\xi z)) \\
    &\eq\N(\eps)^{k_0/2}\sum_{\xi\in(\fr{t}_\lambda\fr{d}^{-1})_+}\N(\fr{t}_\lambda)^{k_0/2}\thin A(\xi\fr{t}_\lambda^{-1}\fr{d})\thin \xi^{-\mbf{k}^\prime/2}\thin e(\tr(\xi z)) \\
    &\eq f(z),
\end{align*}
where $\N(\eps)=1$ since $\eps\in\clo_+^\times$.
\end{proof}

\section{Functional equations}
\label{fqnsection}

In this section, we will prove that the functional equations of $\wh{L}(s,A,\psi)$ imply that $\mbf{f}$ is an eigenfunction of the `Fricke involution'. This section is a direct generalization of the work by Doi and Naganuma in \cite[Section~3]{doinaganuma}. Their proof only deals with the case of parallel weight, level $1$, and class number $1$, and we will modify their proof to fit in the more general case. 

For each Hecke character $\psi$ and each equivalence class $\cl{C}\in\Cl^+(F)$, define a partial $L$-series
\begin{align*}
    L(s,\cl{C},\psi)\eq L(s,A,\cl{C},\psi)\deq\sum_{\fr{a}\in\cl{C}}\psi(\fr{a})\cdot A(\fr{a})\cdot\N(\fr{a})^{-s}.
\end{align*}
Indeed, for $\Re(s)$ sufficiently large,
\begin{align*}
    L(s,A,\psi)\eq\sum_{\cl{C}\in\Cl^+(F)}L(s,\cl{C},\psi).
\end{align*}
Define also the completed $L$-series as
\begin{align}
\label{completedpartiallseries}
    \wh{L}(s,\cl{C},\psi)\deq(2\pi)^{-rs}\; \N(\fr{n}\fr{d}^2)^{s/2}\thin\prod_{j=1}^r\Gamma\bigg(s-\frac{k_j^\prime}{2}+\nu_j(\psi)\bigg)\thin L(s,\cl{C},\psi).
\end{align}

\begin{lemma}
\label{partialfqn}
Let $\cl{C}\in\Cl^+(F)$ be an equivalence class and let $\wt{\cl{C}}=\cl{C}^{-1}\cdot[\fr{n}\fr{d}^2]\in\Cl^+(F)$. Suppose that the $L$-series $L(s)$ satisfies the analytic properties and the functional equations as given in \Cref{mainthm}. Then, for all unramified Hecke characters $\psi$, all the functions $\wh{L}(s,\cl{C},\psi)$ have holomorphic continuation to the whole complex plane which are bounded in vertical strips and satisfy functional equations
\begin{align*}
    \wh{L}(s,\cl{C},\psi)\eq \epsilon \thin i^{-|\mbf{k}|}\thin\psi(\fr{n}\fr{d}^2)\thin\wh{L}(k_0-s,\wt{\cl{C}},\ov{\psi}).
\end{align*}
\end{lemma}
\begin{proof}
Let $G=\Cl^+(F)$ and let $\wh{G}=\Hom(\Cl^+(F),\comp^\times)$ be the dual group. By the theory of characters, $\wh{G}$ is a finite group and the indicator function $\delta_\cl{C}:G\rightarrow\comp$ can be written as
\begin{align*}
    \delta_\cl{C}\eq\sum_{\chi\in\wh{G}}\alpha_\chi\cdot\chi
\end{align*}
for some $\alpha_\chi\in\comp$. Then, for $\Re(s)$ sufficiently large,
\begin{align*}
    L(s,\cl{C},\psi)&\eq\sum_{\fr{a}}\thin\delta_\cl{C}(\fr{a})\cdot\psi(\fr{a})\cdot A(\fr{a})\cdot\N(\fr{a})^{-s} \\
    &\eq \sum_{\chi\in\wh{G}}\alpha_\chi\sum_{\fr{a}}\chi(\fr{a})\cdot\psi(\fr{a})\cdot A(\fr{a})\cdot\N(\fr{a})^{-s} \\
    &\eq\sum_{\chi\in\wh{G}}\alpha_\chi L(s,A,\psi\chi),
\end{align*}
where we also view each $\chi\in\wh{G}$ as a Hecke character of $F$. Hence,
\begin{align*}
    \wh{L}(s,\cl{C},\psi)\eq\sum_{\chi\in\wh{G}}\alpha_\chi \wh{L}(s,A,\psi\chi).
\end{align*}
As each $\wh{L}(s,A,\psi\chi)$ has a holomorphic continuation to the whole complex plane which is bounded in vertical strips, the same holds for $\wh{L}(s,\cl{C},\psi)$. Now, it is easy to check that
\begin{align*}
    \delta_{\wt{\cl{C}}}\eq\sum_{\chi\in\wh{G}}\alpha_\chi\thin\chi(\fr{n}\fr{d}^2)\thin\ov{\chi},
\end{align*}
so for $\Re(s)$ sufficiently large,
\begin{align*}
    \wh{L}(s,\wt{\cl{C}},\ov{\psi})\eq\sum_{\chi\in\wh{G}}\alpha_\chi\thin\chi(\fr{n}\fr{d}^2)\wh{L}(s,A,\ov{\psi\chi}),
\end{align*}
which also holds for all $s\in\comp$ after the holomorphic continuation. By applying the functional equations of $\wh{L}(s,A,\psi\chi)$, one thus obtains that
\begin{align*}
    \wh{L}(s,\cl{C},\psi)\eq\epsilon\thin i^{-|\mbf{k}|}\thin\psi(\fr{n}\fr{d}^2)\sum_{\chi\in\wh{G}}\alpha_\chi\thin\chi(\fr{n}\fr{d}^2)\wh{L}(k_0-s,A,\ov{\psi\chi})\eq\epsilon\thin i^{-|\mbf{k}|}\thin\psi(\fr{n}\fr{d}^2)\thin\wh{L}(k_0-s,\wt{\cl{C}},\ov{\psi}).
\end{align*}
\end{proof}

Recall that for each $1\leq\lambda\leq h$, $\cl{C}_\lambda=[\fr{t}_\lambda^{-1}\fr{d}]$ in $\Cl^+(F)$ and $\wt{\lambda}$ is the unique integer such that $[\fr{t}_\lambda\fr{t}_{\wt{\lambda}}\fr{n}]=[\clo_F]$ in $\Cl^+(F)$. Indeed, for each $1\leq\lambda\leq h$, $\wt{\cl{C}_\lambda}:=\cl{C}_\lambda^{-1}\cdot[\fr{n}\fr{d}^2]=\cl{C}_{\wt{\lambda}}$.

\begin{prop}
\label{fqnprop}
Fix $\lambda$ with $1\leq\lambda\leq h$ and let $\epsilon=\pm1$. Suppose that for any Hecke character $\psi_m$ defined in \Cref{hecke-char-sect}, both $\wh{L}(s,\cl{C}_\lambda,\psi_m)$ and $\wh{L}(s,\cl{C}_{\wt{\lambda}},\ov{\psi_m})$ have holomorphic continuation to the whole complex plane which are bounded in vertical strips and satisfy a functional equation
\begin{align}
\label{fqnequivclass}
    \wh{L}(s,\cl{C}_\lambda,\psi_m)\eq \epsilon \thin i^{-|\mbf{k}|}\thin\psi_m(\fr{n}\fr{d}^2)\thin\wh{L}(k_0-s,\cl{C}_{\wt{\lambda}},\ov{\psi_m}).
\end{align}
Then, $f_\lambda|_\mbf{k}\sm{0&-1\\q_\lambda&0}=\epsilon f_{\wt{\lambda}}$, where $q_\lambda$ is any totally positive generator of $\fr{t}_\lambda\fr{t}_{\wt{\lambda}}\fr{n}$.
\end{prop}

Before proceeding to the proof, we first make another list of notations that we adopt for convenience, which will mostly only appear in this section. Let $y=(y_1,\cdots,y_r)\in(\real_{>0})^r$, $\kappa=(\kappa_1,\cdots,\kappa_r)\in\comp^r$, $s\in\comp$, and $\xi\in F$. We write
\begin{align*}
    iy&\eq(iy_1,\cdots,iy_r),\hspace{0.5cm}\frac{i}{\xi y}\eq\bigg(\frac{i}{\xi^{(1)}y_1},\cdots,\frac{i}{\xi^{(r)}y_r}\bigg) \\
    y^s&\eq\prod_{j=1}^ry_j^s,\hspace{0.5cm} y^\kappa\eq\prod_{j=1}^ry_j^{\kappa_j},\hspace{0.5cm} y^{s+\kappa}\eq y^s\cdot y^\kappa, \hspace{0.5cm}
    \frac{dy}{y}\eq\prod_{j=1}^r\frac{dy_j}{y_j} \\
    \tr(\xi y)&\eq\sum_{j=1}^r \xi^{(j)}y_j, \hspace{0.5cm}\tr(\kappa y)\eq\sum_{j=1}^r\kappa_j y_j
\end{align*}
whenever they make sense. Throughout this section, we also adopt the notations $\psi_m$ and $\nu_m$ defined in \Cref{hecke-char-sect}. Note that this means we have already fixed a basis $\{\eps_1,\cdots,\eps_{r-1}\}$ of~$\clo_+^\times$. 
\\

Now, we go back to the proof. We will first prove some lemmas, which is essentially repeating what Doi and Naganauma did in \cite[Section~3]{doinaganuma}.

\begin{lemma}
\label{gammafunction}
Let $a=(a_j)\in(\real_{>0})^r$ and $b=(b_j)\in\comp^r$. Then, for $\sigma\in\real$ sufficiently large, 
\begin{align*}
    &\frac{1}{2\pi i}\int_{\sigma-i\infty}^{\sigma+i\infty}\prod_{j=1}^r a_j^{-s+b_j}\;\Gamma(s-b_j)\; ds\eq\\
    &\hspace{2cm}r^{-1}\;\Delta\int_{\real^{r-1}}\exp\bigg(-\sum\limits_{j=1}^r \big(a_j\prod\limits_{l=1}^{r-1}(\eps_l^{(j)})^{t_l}\big)\bigg)\;\prod_{j=1}^r\prod_{l=1}^{r-1}(\eps_l^{(j)})^{-t_lb_j}\;dt_1\cdots dt_{r-1},
\end{align*}
where
\begin{align*}
    \Delta\eq\begin{vmatrix}1&\log(\eps_1^{(1)})&\cdots&\log(\eps_{r-1}^{(1)})\\ 1&\log(\eps_1^{(2)})&\cdots&\log(\eps_{r-1}^{(2)})\\
    \vdots&\vdots & &\vdots \\ 1&\log(\eps_1^{(r)})&\cdots&\log(\eps_{r-1}^{(r)}) \end{vmatrix}.
\end{align*}
\end{lemma}
\begin{proof}
First, note that
\begin{align*}
    \prod_{j=1}^r a_j^{-s+b_j}\;\Gamma(s-b_j)\eq\prod_{j=1}^r\int_{\real_{>0}}e^{-a_jy_j}\; y_j^{s-b_j}\;\frac{y_j}{y_j} \eq\int_{(\real_{>0})^r}e^{-\tr(ay)}\;y^{s-b}\;\frac{dy}{y}
\end{align*}
We do a change of variables $(y_1,\cdots,y_r)\mapsto(y_0,t_1,\cdots,t_{r-1})$ as the following
\begin{align*}
    y_j\eq y_0\thin\prod_{l=1}^{r-1}(\eps_l^{(j)})^{t_l}
\end{align*}
where $y_0\in\real_{>0}$ and $t_l\in\real$. It is easy to check that
\begin{align*}
    y_1y_2\cdots y_r&\eq y_0^r \\
    dy_1dy_2\cdots dy_r&\eq \Delta\;y_0^{r-1}\;dy_0\;dt_1\cdots dt_{r-1}.
\end{align*}
Hence,
\begin{align*}
    \prod_{j=1}^r a_j^{-s+b_j}\;\Gamma(s-b_j)&\eq\int_{(\real_{>0})^r}e^{-\tr(ay)}\;y^{s-b}\;\frac{dy}{y} \\
    &\eq\Delta\int_{\real_{>0}}\bigg(\int_{\real^{r-1}}\exp\bigg(-\sum\limits_{j=1}^r a_jy_j\bigg)y^{-b}\;dt_1\cdots dt_{r-1}\bigg)y_0^{rs}\;\frac{dy_0}{y_0} \\
    &\eq r^{-1}\;\Delta\int_{\real_{>0}}\bigg(\int_{\real^{r-1}}\exp\bigg(-\sum\limits_{j=1}^r a_jy_j\bigg)y^{-b}\;dt_1\cdots dt_{r-1}\bigg)u^s\;\frac{du}{u}
\end{align*}
where we still keep the variable $y_j=y_0\prod_{l=1}^{r-1}(\eps_l^{(j)})^{t_l}$ for convenience and do another change of variables $u=y_0^r$. By the inverse Mellin transform, we may pick $\sigma\in\real$ sufficiently large such that
\begin{align*}
    &\frac{1}{2\pi i}\int_{\sigma-i\infty}^{\sigma+i\infty}\prod_{j=1}^r a_j^{-s+b_j}\thin\Gamma(s-b_j)\thin u^{-s}\thin ds\eq r^{-1}\thin\Delta\int_{\real^{r-1}}\exp\bigg(-\sum\limits_{j=1}^r a_jy_j\bigg)y^{-b}\thin dt_1\cdots dt_{r-1}
\end{align*}
where we still keep the variable $y_j=u^{1/r}\prod_{l=1}^{r-1}(\eps_l^{(j)})^{t_l}$ for convenience. Plugging in $u=1$ on both sides then gives the result.
\end{proof}

Before stating the corollary, we write
\begin{align*}
    \eps(t)&\eq \prod_{l=1}^{r-1}\eps_l^{t_l} \\
    \eps(t)^c&\eq\prod_{j=1}^r\prod_{l=1}^{r-1}(\eps_l^{(j)})^{t_lc_j} \\
    \tr(a\eps(t))&\eq\sum_{j=1}^r \bigg(a_j\prod_{l=1}^{r-1}(\eps_l^{(j)})^{t_l}\bigg)
\end{align*}
for $t\in\real^{r-1}$, $c\in\real^r$, and $a\in\comp^r$ whenever they make sense.

\begin{corol}
\label{gammainversemellin-corol}
Fix $m\in\inte^{r-1}$ and let $\{\nu_{m,j}\}_{j=1}^r$ be defined as in \Cref{hecke-char-sect}. Let $\mbf{k}^\prime$ be as before. Let $a=(a_j)\in(\real_{>0})^r$. Then, 
\begin{equation}
\begin{aligned}
\label{gammainversemellin}
    &\frac{1}{2\pi i}\int_{\sigma-i\infty}^{\sigma+i\infty}\prod_{j=1}^r a_j^{-s+(k_j^\prime/2)-\nu_{m,j}}\;\Gamma\bigg(s-\frac{k_j^\prime}{2}+\nu_{m,j}\bigg)\thin ds\eq \\
    &\hspace{4cm}r^{-1}\;\Delta\int_{\real^{r-1}}e^{-\tr(a\eps(t))}\;\eps(t)^{-\mbf{k}^\prime/2}\thin e\bigg(\sum_{l=1}^{r-1}m_lt_l\bigg)\thin dt_1\cdots dt_{r-1}.
\end{aligned}
\end{equation}
\end{corol}
\begin{proof}
Apply \Cref{gammafunction} with $b_j=\frac{k_j^\prime}{2}-\nu_{m,j}$ and use that $\eps(t)^{\nu_m}=e\big(\sum_{l=1}^{r-1}m_lt_l\big)$.
\end{proof}

\begin{lemma}
\label{phibeta}
Let $\beta\in F_+$ be a totally positive element. Then,
\begin{equation}
\begin{aligned}
\label{phibetaeqn}
    &\sum_{\eps\in\clo_+^\times}(\eps\beta)^{-\mbf{k}^\prime/2}\;e^{-2\pi\cdot\tr(\eps\beta y)}\eq \\
    &\hspace{1cm}C\sum_{m\in\inte^{r-1}}\int_{\sigma-i\infty}^{\sigma+i\infty}(2\pi)^{-rs}\thin\N(\beta)^{-s}\thin\psi_m((\beta))\thin y^{-s+(\mbf{k}^\prime/2)-\nu_m}\prod_{j=1}^r\Gamma\bigg(s-\frac{k_j^\prime}{2}+\nu_{m,j}\bigg)\thin ds
\end{aligned}
\end{equation}
for $y\in(\real_{>0})^r$, where $C=\frac{r\Delta^{-1}(2\pi)^{|\mbf{k}^\prime|/2}}{2\pi i}$.
\end{lemma}
\begin{proof}
Let
\begin{align*}
    \vphi_\beta(y)\deq\sum_{\eps\in\clo_+^\times}(\eps\beta)^{-\mbf{k}^\prime/2}\;e^{-2\pi\cdot\tr(\eps\beta y)}\eq\beta^{-\mbf{k}^\prime/2}\sum_{m\in\inte^{r-1}}\eps(m)^{-\mbf{k}^\prime/2}\;e^{-2\pi\cdot\tr(\eps(m)\beta y)}.
\end{align*}
Write
\begin{align*}
    S(t)\deq \eps(t)^{-\mbf{k}^\prime/2}\;e^{-2\pi\cdot\tr(\eps(t)\beta y)}\hspace{0.5cm}\text{and}\hspace{0.5cm}\wh{S}(x)\deq\int_{\real^{r-1}}S(t)\; e\bigg(-\sum_{l=1}^{r-1}x_lt_l\bigg)\thin dt
\end{align*}
for $t,x\in\real^{r-1}$. Then, applying \Cref{gammainversemellin-corol} with $a=(2\pi\beta^{(j)}y_j)_j\in(\real_{>0})^r$, we have
\begin{align*}
    \wh{S}(-m)&\eq\int_{\real^{r-1}}\eps(t)^{-\mbf{k}^\prime/2}\;e^{-2\pi\cdot\tr(\eps(t)\beta y)}\; e\bigg(\sum_{l=1}^{r-1}m_lt_l\bigg)\thin dt_1\cdots dt_{r-1} \\
    &\;\stackrel{(\ref{gammainversemellin})}{=}\;\frac{r\Delta^{-1}}{2\pi i}\int_{\sigma-i\infty}^{\sigma+i\infty}\prod_{j=1}^r(2\pi\beta^{(j)}y_j)^{-s+(k_j^\prime/2)-\nu_{m,j}}\;\Gamma\bigg(s-\frac{k_j^\prime}{2}+\nu_{m,j}\bigg)\thin ds
\end{align*}
for $m\in\inte^{r-1}$. By the Poisson summation formula,
\begin{align*}
    \vphi_\beta(y)&\eq\beta^{-\mbf{k}^\prime/2}\sum_{m\in\inte^{r-1}}S(m)\eq\beta^{-\mbf{k}^\prime/2}\sum_{m\in\inte^{r-1}}\wh{S}(m) \\
    &\eq\frac{r\Delta^{-1}\beta^{-\mbf{k}^\prime/2}}{2\pi i}\sum_{m\in\inte^{r-1}}\int_{\sigma-i\infty}^{\sigma+i\infty}\prod_{j=1}^r(2\pi\beta^{(j)}y_j)^{-s+(k_j^\prime/2)-\nu_{m,j}}\;\Gamma\bigg(s-\frac{k_j^\prime}{2}+\nu_{m,j}\bigg)\thin ds \\
    &\eq C\sum_{m\in\inte^{r-1}}\int_{\sigma-i\infty}^{\sigma+i\infty}(2\pi)^{-rs}\;\N(\beta)^{-s}\;\psi_m((\beta))\;y^{-s+(\mbf{k}^\prime/2)-\nu_m}\prod_{j=1}^r\Gamma\bigg(s-\frac{k_j^\prime}{2}+\nu_{m,j}\bigg)\thin ds,
\end{align*}
where $C=\frac{r\Delta^{-1}(2\pi)^{|\mbf{k}^\prime|/2}}{2\pi i}$.
\end{proof}

We are now ready to prove \Cref{fqnprop}. Define $\vphi_\beta(y)$ as in the proof of \Cref{phibeta}. Applying \Cref{phibeta} with $\beta=\xi\in(\fr{t}_\lambda\fr{d}^{-1})_+$ and letting $\fr{a}=\xi\fr{t}_\lambda^{-1}\fr{d}\in\cl{C}_\lambda$, we obtain that
\begin{align*}
    \vphi_{\xi}(y)&\;\stackrel{(\ref{phibetaeqn})}{=}\;C\sum_{m\in\inte^{r-1}}\int_{\sigma-i\infty}^{\sigma+i\infty}(2\pi)^{-rs}\;\N(\xi)^{-s}\;\psi_m((\xi))\;y^{-s+(\mbf{k}^\prime/2)-\nu_m}\prod_{j=1}^r\Gamma\bigg(s-\frac{k_j^\prime}{2}+\nu_{m,j}\bigg)\thin ds \\
    &\eq C\sum_{m\in\inte^{r-1}}\int_{\sigma-i\infty}^{\sigma+i\infty}(2\pi)^{-rs}\;\N(\fr{a}\fr{t}_\lambda\fr{d}^{-1})^{-s}\;\psi_m(\fr{a}\fr{t}_\lambda\fr{d}^{-1}) \\
    &\hspace{5cm}\times y^{-s+(\mbf{k}^\prime/2)-\nu_m}\prod_{j=1}^r\Gamma\bigg(s-\frac{k_j^\prime}{2}+\nu_{m,j}\bigg)\thin ds.
\end{align*}
Now, 
\begin{align*}
    f_\lambda(iy)&\eq\sum_{\xi\in(\fr{t}_\lambda\fr{d}^{-1})_+}\N(\fr{t}_\lambda)^{k_0/2}A(\xi\fr{t}_\lambda^{-1}\fr{d})\thin\xi^{-\mbf{k}^\prime/2}\thin e^{-2\pi \cdot\tr(\xi y)} \\
    &\eq\N(\fr{t}_\lambda)^{k_0/2}\sum_{\xi\in(\fr{t}_\lambda\fr{d}^{-1})_+/\clo_+^\times}\sum_{\eps\in\clo_+^\times}A(\xi\fr{t}_\lambda^{-1}\fr{d})\thin(\xi\eps)^{-\mbf{k}^\prime/2}\thin e^{-2\pi \cdot\tr(\xi\eps y)} \\
    &\eq\N(\fr{t}_\lambda)^{k_0/2}\sum_{\xi\in(\fr{t}_\lambda\fr{d}^{-1})_+/\clo_+^\times}A(\xi\fr{t}_\lambda^{-1}\fr{d})\thin\vphi_\xi(y).
\end{align*}
Applying our computations for $\vphi_\xi(y)$ and letting $\fr{a}=\xi\fr{t}_\lambda^{-1}\fr{d}$, we obtain that
\begin{align}
    f_\lambda(iy)&\eq C\thin\N(\fr{t}_\lambda)^{k_0/2}\sum_{\fr{a}\in\cl{C}_\lambda}\sum_{m\in\inte^{r-1}}\int_{\sigma-i\infty}^{\sigma+i\infty}(2\pi)^{-rs}\;\;\N(\fr{a}\fr{t}_\lambda\fr{d}^{-1})^{-s}\;\psi_m(\fr{a}\fr{t}_\lambda\fr{d}^{-1})\;A(\fr{a}) \nonumber \\
    &\hspace{7cm}\times y^{-s+(\mbf{k}^\prime/2)-\nu_m}\;\prod_{j=1}^r\Gamma\bigg(s-\frac{k_j^\prime}{2}+\nu_{m,j}\bigg)\thin ds \nonumber \\
    &\;\stackrel{(\ref{completedpartiallseries})}{=}\; C\sum_{m\in\inte^{r-1}}\int_{\sigma-i\infty}^{\sigma+i\infty}\psi_{m}(\fr{t}_\lambda\fr{d}^{-1})\;\N(\fr{n})^{-s/2}\;\N(\fr{t}_\lambda)^{(k_0/2)-s}\;y^{-s+(\mbf{k}^\prime/2)-\nu_m}\;\wh{L}(s,\cl{C}_\lambda,\psi_m)\; ds \label{flambdaiy}
\end{align}
By the functional equations,
\begin{align*}
    &f_\lambda(iy)\;\stackrel{(\ref{fqnequivclass})}{=}\; \\
    &\hspace{0.5cm}\epsilon\thin i^{-|\mbf{k}|}\thin C\sum_{m\in\inte^{r-1}}\int_{\sigma-i\infty}^{\sigma+i\infty}\psi_m(\fr{n}\fr{t}_\lambda\fr{d})\;\N(\fr{n})^{-s/2}\;\N(\fr{t}_\lambda)^{(k_0/2)-s}\; y^{-s+(\mbf{k}^\prime/2)-\nu_m}\;\wh{L}(k_0-s,\cl{C}_{\wt{\lambda}},\psi_{-m})\; ds.
\end{align*}
By moving the line of the integral (which is possible under the assumed analytic properties of all $\wh{L}(s,\cl{C}_\lambda,\psi_m)$), we obtain that
\begin{align*}
    &f_\lambda(iy)\\
    \eq&\epsilon\thin C\thin (iy)^{-\mbf{k}}\sum_{m\in\inte^{r-1}}\int_{\sigma-i\infty}^{\sigma+i\infty}\psi_m(\fr{n}\fr{t}_\lambda\fr{d})\;\N(\fr{n})^{(s-k_0)/2}\;\N(\fr{t}_\lambda)^{s-(k_0/2)}\;y^{s-(\mbf{k}^\prime/2)-\nu_m}\;\wh{L}(s,\cl{C}_{\wt{\lambda}},\psi_{-m})\; ds \\
    \eq&\epsilon\thin C\thin (iy)^{-\mbf{k}}\sum_{m\in\inte^{r-1}}\int_{\sigma-i\infty}^{\sigma+i\infty}\psi_{-m}(\fr{n}\fr{t}_\lambda\fr{d})\;\N(\fr{n})^{(s-k_0)/2}\;\N(\fr{t}_\lambda)^{s-(k_0/2)}\;y^{s-(\mbf{k}^\prime/2)+\nu_m}\;\wh{L}(s,\cl{C}_{\wt{\lambda}},\psi_m)\; ds
\end{align*}
On the other hand, applying the computations for $f_\lambda(iy)$ to $f_{\wt{\lambda}}(i/(q_\lambda y))$, we obtain that
\begin{align*}
    &\epsilon\thin q_\lambda^{-\mbf{k}/2}\thin (iy)^{-\mbf{k}}\thin f_{\wt{\lambda}}\bigg(\frac{i}{q_\lambda y}\bigg)\\
    \;\stackrel{(\ref{flambdaiy})}{=}\;&\epsilon\thin C\thin q_\lambda^{-\mbf{k}/2}\thin(iy)^{-\mbf{k}}\sum_{m\in\inte^{r-1}}\int_{\sigma-i\infty}^{\sigma+i\infty}\psi_{m}(\fr{t}_{\wt{\lambda}}\fr{d}^{-1})\;\N(\fr{n})^{-s/2}\;\N(\fr{t}_{\wt{\lambda}})^{(k_0/2)-s} \\
    &\hspace{8cm}\times(1/q_\lambda y)^{-s+(\mbf{k}^\prime/2)-\nu_m}\;\wh{L}(s,\cl{C}_{\wt{\lambda}},\psi_m)\; ds \\
    \eq&\epsilon\thin C\thin(iy)^{-\mbf{k}}\sum_{m\in\inte^{r-1}}\int_{\sigma-i\infty}^{\sigma+i\infty}\psi_{m}(q_\lambda^{-1}\fr{t}_{\wt{\lambda}}\fr{d}^{-1})\;\N(\fr{n})^{-s/2}\;\N(\fr{t}_{\wt{\lambda}})^{(k_0/2)-s}\;\N(q_\lambda)^{s-(k_0/2)} \\
    &\hspace{8cm}\times y^{s-(\mbf{k}^\prime/2)+\nu_m}\;\wh{L}(s,\cl{C}_{\wt{\lambda}},\psi_m)\; ds \\
    \eq&\epsilon\thin C\thin(iy)^{-\mbf{k}}\sum_{m\in\inte^{r-1}}\int_{\sigma-i\infty}^{\sigma+i\infty}\psi_{-m}(\fr{n}\fr{t}_\lambda\fr{d})\;\N(\fr{n})^{(s-k_0)/2}\;\N(\fr{t}_\lambda)^{s-(k_0/2)}\;y^{s-(\mbf{k}^\prime/2)+\nu_m}\;\wh{L}(s,\cl{C}_{\wt{\lambda}},\psi_m)\; ds,
\end{align*}
where we use that $(q_\lambda)=\fr{t}_\lambda\fr{t}_{\wt{\lambda}}\fr{n}$. Hence, by comparing the two equations, we obtain that
\begin{align*}
    f_\lambda(iy)\eq\epsilon\thin q_\lambda^{-\mbf{k}/2}\thin (iy)^{-\mbf{k}}\thin f_{\wt{\lambda}}\bigg(\frac{i}{q_\lambda y}\bigg)
\end{align*}
for $y\in(\real_{>0})^r$. By analytic continuation, this implies that
\begin{align*}
    f_\lambda(z)\eq \epsilon\thin q_\lambda^{-\mbf{k}/2}\thin z^{-\mbf{k}}\thin f_{\wt{\lambda}}\bigg(-\frac{1}{q_\lambda z}\bigg)\eq \epsilon f_{\wt{\lambda}}|_\mbf{k}\m{0&-1\\q_\lambda&0}.
\end{align*}
As $\mbf{k}\in(2\inte^+)$, this is equivalent to $f_\lambda|_\mbf{k}\sm{0&-1\\q_\lambda&0}=\epsilon f_{\wt{\lambda}}$.

\begin{corol}
\label{fqncorol}
If the $L$-series $L(s)$ satisfies the analytic properties and the functional equations as given in \Cref{mainthm}, then 
\begin{enumerate}[label={\rm(\arabic*)}]
    \item $f_\lambda|_\mbf{k}W_{q_\lambda}=\epsilon f_{\wt{\lambda}}$ for all $1\leq \lambda\leq h$, where $W_{q_\lambda}:=\sm{0&-1\\q_\lambda&0}$ and $q_\lambda$ is any totally positive generator of $\fr{t}_\lambda\fr{t}_{\wt{\lambda}}\fr{n}$;
    \item $f_\lambda|_\mbf{k}A_\beta=f_\lambda$ for all $\beta\in\fr{t}_\lambda\fr{n}$, where $A_\beta:=\sm{1&0\\\beta&1}$.
\end{enumerate}
\end{corol}
\begin{proof}
(1) follows from \Cref{partialfqn} and \Cref{fqnprop}. For (2), note that for any $\beta\in\fr{t}_\lambda\fr{n}$, $\beta/q_\lambda\in\fr{t}_{\wt{\lambda}}^{-1}$, and that $A_\beta=W_{q_\lambda}^{-1}T^{-\beta/q_\lambda}W_{q_\lambda}$. Hence, 
\begin{align*}
    f_\lambda|_\mbf{k}A_\beta\eq f_\lambda|_\mbf{k}W_{q_\lambda}^{-1}T^{-\beta/q_\lambda}W_{q_\lambda}\eq\epsilon f_{\wt{\lambda}}|_\mbf{k}T^{-\beta/q_\lambda}W_{q_\lambda}\eq\epsilon f_{\wt{\lambda}}|_\mbf{k}W_{q_\lambda}\eq f_\lambda
\end{align*}
for all $\beta\in\fr{t}_\lambda\fr{n}$.
\end{proof}

\section{Generating sets of congruence subgroups}
\label{generating-sect}
In this section, we will recall the main theorem in \cite{vaserstein} and prove some important results on the generating sets of certain congruence subgroups, which will be used for the proof of \Cref{mainthm}. Throughout this section, we let $K$ denote a general number field to distinguish it from $F$, which we reserve for our totally real field of degree $r>1$.

We first recall the definition of Dedekind domains of arithmetic type from \cite{dedearith}, which was used in \cite{vaserstein}. Let $K$ be a number field and $\Sigma$ be a finite set of places of $K$ containing all the archimedean ones. Define
\begin{align*}
    \clo_{K,\Sigma}\eq\{\alpha\in K\mid \val_v(\alpha)\geq 0\text{ for all }v\notin \Sigma\}.
\end{align*}
A ring is called a \textit{Dedekind domain of arithmetic type} if it is of the form $\clo_{K,\Sigma}$ for some number field~$K$ and some finite set $\Sigma$ of places of $K$ containing all the archimedean ones.

Let $A=\clo_{K,\Sigma}$ be a Dedekind domain of arithmetic type for some $K$ and $\Sigma$ and let $I_1,I_2$ be two nonzero ideals of $A$. Define
\begin{align*}
    G(I_1,I_2)\deq\bigg\{\m{a&b\\c&d}\in\SL_2(A)\;\bigg|\; a-1,d-1\in I_1I_2,\;b\in I_1,\;c\in I_2\bigg\}
\end{align*}
and define $E(I_1,I_2)$ to be the subgroup of $\SL_2(A)$ generated by matrices of the form $\sm{1&b\\0&1}$ and $\sm{1&0\\c&1}$ for $b\in I_1$ and $c\in I_2$. Indeed, $E(I_1,I_2)\subseteq G(I_1,I_2)$ and $G(I_1,I_2)$ is a finite index subgroup of $\SL_2(A)$.

\begin{theorem}[Vaserstein \cite{vaserstein}]
\label{vaserstein}Let $A=\clo_{K,\Sigma}$ be a Dedekind domain of arithmetic type for some $K$ and $\Sigma$. Suppose that $A$ has infinitely many units, or equivalently, $\Sigma$ contains at least two places. Then, for any nonzero ideals $I_1,I_2$ of $A$, $E(I_1,I_2)$ is a normal and finite index subgroup of $G(I_1,I_2)$, and $E(A,A)=G(A,A)=\SL_2(A)$. Moreover, if $A$ is not an order in a totally imaginary number field, then $E(I_1,I_2)=G(I_1,I_2)$.
\end{theorem}
\begin{rmk}
There is a small gap in the proof of \Cref{vaserstein} as the validity of \cite[Lemma~3]{vaserstein} was later called into question, but the proof was rewritten in \cite{liehl} without further controversy. One should see \cite{liehl} for more details.
\end{rmk}
\begin{rmk}
A straightforward corollary is that if $K$ is a real quadratic field, then $\SL_2(\clo_K)$ is generated by matrices of the form $\sm{1&b\\0&1}$ and $\sm{1&0\\c&1}$ for \mbox{$b,c\in\clo_K$}. As mentioned in the introduction, together with the work by Doi and Nagaunuma in \cite{doinaganuma}, this yields a converse theorem for Hilbert modular forms of level $1$, subject to the class number $1$ condition.
\end{rmk}

Now, we would like to apply \Cref{vaserstein} to our setting. We will be writing results in terms of general number fields so that they may be applied to more general cases. Let $K$ be a general number field. Define
\begin{align*}
    \GL_2^+(K)\deq\{\gamma\in\GL_2(K)\mid\sigma(\det\gamma)>0\text{ for all }\sigma:K\hookrightarrow\real\}.
\end{align*}
For nonzero ideals $\fr{t}$ and $\fr{m}$ of $\clo_K$, define
\begin{align*}
    \SG_1(\fr{m})&\deq\bigg\{\m{a&b\\c&d}\in\SL_2(\clo_K)\;\bigg|\; \m{a&b\\c&d}\equiv\m{1&*\\0&1}\pmod{\fr{m}}\bigg\} \\
    \Gamma_1(\fr{t},\fr{m})&\deq\bigg\{\m{a&b\\c&d}\in\GL_2(K)\;\bigg|\; a\in\clo_K,\;b\in\fr{t}^{-1},\;c\in\fr{t}\fr{m},\;d-1\in\fr{m},\;ad-bc\in\clo_K^\times\bigg\} \\
    \Gamma_1^+(\fr{t},\fr{m})&\deq\bigg\{\m{a&b\\c&d}\in\GL_2(K)\;\bigg|\; a\in\clo_K,\;b\in\fr{t}^{-1},\;c\in\fr{t}\fr{m},\;d-1\in\fr{m},\;ad-bc\in(\clo_K^\times)_+\bigg\}\\
    &\;\eq\Gamma_1(\fr{t},\fr{m})\cap\GL_2^+(K).
\end{align*}
Here $(\clo_K^\times)_+$ denotes the set of totally positive units, i.e., $\eps\in\clo_K^\times$ such that $\sigma(\eps)>0$ for all $\sigma:K\hookrightarrow\real$.

\begin{lemma}
\label{sgamma1-generator}
Let $K$ be a number field of degree greater than $1$ and having at least one real place. Let $\fr{m}$ be a nonzero ideal of $\clo_K$. Then, $\SG_1(\fr{m})$ is generated by matrices of the form $\sm{1&b\\0&1}$ and $\sm{1&0\\c&1}$ for $b\in\clo_K$ and $c\in\fr{m}$.
\end{lemma}
\begin{proof}
By the assumption on $K$, $K$ has at least two archimedean places and $\clo_K$ is not an order in a totally imaginary number field. Hence, we may apply \Cref{vaserstein} to $A=\clo_K$, $I_1=\clo_K$, and $I_2=\fr{m}$ to obtain the result.
\end{proof}

\begin{prop}
\label{gamma1-generator}
Let $K$ be a number field of degree greater than $1$ and having at least one real place. Let $\fr{t}$ and $\fr{m}$ be nonzero ideals of $\clo_K$. Then, $\Gamma_1(\fr{t},\fr{tm})$ (resp. $\Gamma_1^+(\fr{t},\fr{tm})$) is contained in the group generated by matrices of the form $\sm{1&b\\0&1}$, $\sm{1&0\\c&1}$, and $\sm{\eps&0\\0&1}$ for $b\in\fr{t}^{-1}$, $c\in\fr{tm}$, and $\eps\in\clo_K^\times$ (resp. $\eps\in(\clo_K^\times)_+$).
\end{prop}
\begin{proof}
Let $\gamma\in\Gamma_1(\fr{t},\fr{tm})$ or $\Gamma_1^+(\fr{t},\fr{tm})$. If we write $\eps=\det\gamma$, then we may replace $\gamma$ by $\sm{\eps^{-1}&0\\0&1}\gamma$ to assume that $\det\gamma=1$. In this way, it suffices to show that any $\gamma\in\Gamma_1(\fr{t},\fr{tm})$ with $\det\gamma=1$ is generated by matrices of the form $\sm{1&\alpha\\0&1}$ and $\sm{1&0\\\beta&1}$ for $\alpha\in\fr{t}^{-1}$ and $\beta\in\fr{tm}$.

Write $\gamma=\sm{a&b\\c&d}$. Note that the ideals $(a)$ and $\fr{t}$ must be coprime, as otherwise their common prime factor would divide $(ad-bc)$ (since $\fr{tm}|(bc)$), contradicting that $ad-bc=1$. Now, let $S$ be a set of representatives of $\clo_K/\fr{t}$. Since $(a)$ and $\fr{t}$ are coprime, $\{an+1|n\in S\}$ is also a set of representatives of $\clo_K/\fr{t}$. In particular, there exists $n\in\clo_K$ such that $an+1\in\fr{t}$. Note that
\begin{align*}
    \gamma\m{1&bn\\0&1}\eq\m{a&b(an+1)\\c&bcn+d}.
\end{align*}
Here, we have
\begin{enumerate}[label=(\arabic*)]
    \item $a\in\clo_K$;
    \item $b(an+1)\in\clo_K$ since $b\in\fr{t}^{-1}$ and $an+1\in\fr{t}$;
    \item $c\in\fr{t}^2\fr{m}\subseteq\fr{tm}$;
    \item $bcn+d-1\in\fr{tm}$ since $b\in\fr{t}^{-1}$, $c\in\fr{t}^2\fr{m}$, and $d-1\in\fr{tm}$.
\end{enumerate}
As $ad-bc=1$, we also have $a-1\in\fr{tm}$. Hence, $\gamma\sm{1&bn\\0&1}\in\SG_1(\fr{tm})$. By \Cref{sgamma1-generator}, $\SG_1(\fr{tm})$ is generated by matrices of the form $\sm{1&\alpha\\0&1}$ and $\sm{1&0\\\beta&1}$ for $\alpha\in\clo_K$ and $\beta\in\fr{tm}$, so the result follows.
\end{proof}

We will also prove a result for $\Gamma_0$ as we will need it later, the proof of which does not rely on \Cref{vaserstein} though. Define
\begin{align*}
    \SG_0(\fr{m})&\deq\bigg\{\m{a&b\\c&d}\in\SL_2(\clo_K)\;\bigg|\; c\in\fr{m}\bigg\} \\
    \Gamma_0(\fr{t},\fr{m})&\deq\bigg\{\m{a&b\\c&d}\in\GL_2(K)\;\bigg|\; a\in\clo_K,\;b\in\fr{t}^{-1},\;c\in\fr{t}\fr{m},\;d\in\clo_K,\;ad-bc\in\clo_K^\times\bigg\} \\
    \Gamma_0^+(\fr{t},\fr{m})&\deq\bigg\{\m{a&b\\c&d}\in\GL_2(K)\;\bigg|\; a\in\clo_K,\;b\in\fr{t}^{-1},\;c\in\fr{t}\fr{m},\;d\in\clo_K,\;ad-bc\in(\clo_K^\times)_+\bigg\}\\
    &\;\eq\Gamma_0(\fr{t},\fr{m})\cap\GL_2^+(K).
\end{align*}

\begin{prop}
\label{gamma0-generator}
Let $K$ be a number field. Let $\fr{t}$ and $\fr{m}$ be nonzero ideals of $\clo_K$ such that $\fr{t}$ is a \textbf{prime} ideal. Then, $\Gamma_0(\fr{t},\fr{m})$ (resp. $\Gamma_0^+(\fr{t},\fr{m})$) is generated by $\SG_0(\fr{tm})$ and matrices of the form $\sm{1&b\\0&1}$ and $\sm{\eps&0\\0&1}$ for $b\in\fr{t}^{-1}$ and $\eps\in\clo_K^\times$ (resp. $\eps\in(\clo_K^\times)_+$).
\end{prop}
\begin{proof}
Similarly as in the proof of \Cref{gamma1-generator}, it suffices to show that any $\gamma\in\Gamma_0(\fr{t},\fr{m})$ with $\det\gamma=1$ is generated by $\SG_0(\fr{tm})$ and matrices of the form $\sm{1&\alpha\\0&1}$ for $\alpha\in\fr{t}^{-1}$.

First, suppose that $a\in\fr{t}$. Let $\fr{q}\neq\fr{t}$ be another prime ideal with $[\fr{q}]=[\fr{t}]$ in the class group $\Cl(K)$ of $K$ (such $\fr{q}$ exists by Chebotarev's density theorem). Then, there exists $l\in K$ such that $(l)\fr{t}=\fr{q}$ (and this also implies that $l\in\fr{t}^{-1}$). Suppose for the sake of contradiction that $a+cl\in\fr{t}$. Then, $cl\in\fr{t}$ since $a\in\fr{t}$. Note that $(cl)=(c)\fr{t}^{-1}\fr{q}$, where $(c)\fr{t}^{-1}$ is an integral ideal since $c\in\fr{tm}$. As $\fr{t}|(cl)$ and $(\fr{t},\fr{q})=1$, this implies that $\fr{t}|(c)\fr{t}^{-1}$. Thus, $\fr{t}|(bc)$ since $b\in\fr{t}^{-1}$ and hence $\fr{t}|(ad-bc)$ since $a\in\fr{t}$ by the assumption, which leads to contradiction as $\fr{t}$ is a prime ideal while $ad-bc=\det\gamma=1$. Hence, $a+cl\notin\fr{t}$. This means that we may replace $\gamma$ by $\sm{1&l\\0&1}\gamma$ to assume that $a\notin\fr{t}$.

Now, $(a)$ and $\fr{t}$ are coprime, so similarly as in the proof of \Cref{gamma1-generator}, there exists $n\in\clo_K$ such that $an+1\in\fr{t}$. Then,
\begin{align*}
    \gamma\m{1&bn\\0&1}\eq\m{a&b(an+1)\\c&bcn+d}
\end{align*}
with
\begin{enumerate}[label=(\arabic*)]
    \item $a\in\clo_K$;
    \item $b(an+1)\in\clo_K$ since $b\in\fr{t}^{-1}$ and $an+1\in\fr{t}$;
    \item $c\in\fr{tm}$;
    \item $bcn+d\in\clo_K$ since $b\in\fr{t}^{-1}$, $c\in\fr{tm}$, and $d\in\clo_K$.
\end{enumerate}
Hence, $\gamma\sm{1&bn\\0&1}\in\SG_0(\fr{tm})$. The result then follows.
\end{proof}

\section{Invariance of $f_\lambda$ under $\Gamma_1^+(\fr{t}_\lambda,\fr{m}\fr{n})$}
\label{gamma1-inv}

We will now give a proof of the first part of \Cref{mainthm}, i.e., recovering a Hilbert modular form assuming only the `unramified' functional equations.

\begin{prop}
\label{f_lambda-gamma1-inv}
Let $\{\fr{t}_1,\cdots,\fr{t}_h\}$ be a set of representatives of $\Cl^+(F)$. If the $L$-series $L(s)$ satisfies the analytic properties and the functional equations as given in \Cref{mainthm}, then for each \mbox{$1\leq\lambda\leq h$}, the associated holomorphic function $f_\lambda$ satisfies that $f_\lambda|_\mbf{k}\gamma=f_\lambda$ for all $\gamma\in\Gamma_1^+(\fr{t}_\lambda,\fr{t}_\lambda\fr{n})$.
\end{prop}
\begin{proof}
By \Cref{basic} and \Cref{fqncorol}, we have that for all $1\leq\lambda\leq h$,
\begin{enumerate}[label=\rm{(\alph*)}]
    \item $f_\lambda|_\mbf{k}T^\alpha=f_\lambda$ for all $\alpha\in\fr{t}_\lambda^{-1}$, where $T^\alpha=\sm{1&\alpha\\0&1}$;
    \item $f_\lambda|_\mbf{k}E_\eps=f_\lambda$ for all $\eps\in\clo_+^\times$, where $E_\eps=\sm{\eps&0\\0&1}$;
    \item $f_\lambda|_\mbf{k}A_\beta=f_\lambda$ for all $\beta\in\fr{t}_\lambda\fr{n}$, where $A_\beta=\sm{1&0\\\beta&1}$. 
\end{enumerate}
As $F$ is a totally real field of degree $r>1$, it follows from \Cref{gamma1-generator} that $f_\lambda$ is invariant under $\Gamma_1^+(\fr{t}_\lambda,\fr{t}_\lambda\fr{n})$. 
\end{proof}

\begin{corol}
\label{f-gamma1-inv}
Let $\{\fr{t}_1,\cdots,\fr{t}_h\}$ be a set of representatives of $\Cl^+(F)$ and let $\fr{m}=\fr{t}_1\cdots\fr{t}_h$. If the $L$-series $L(s)$ satisfies the analytic properties and the functional equations as given in \Cref{mainthm}, then the associated $h$-tuple $\mbf{f}=(f_1,\cdots,f_h)$ is a Hilbert modular form of weight $\mbf{k}$ and level $K_1(\fr{mn})$.
\end{corol}
\begin{proof}
By \Cref{f_lambda-gamma1-inv}, each $f_\lambda$ is a classical Hilbert modular form of weight $\mbf{k}$ and level $\Gamma_1^+(\fr{t}_\lambda,\fr{t}_\lambda\fr{n})$, so in particular, it is also of level $\Gamma_1^+(\fr{t}_\lambda,\fr{m}\fr{n})$. The result then follows.
\end{proof}

The first part of \Cref{mainthm} now simply follows from \Cref{f-gamma1-inv} and \Cref{fhmf}. 

\section{Euler product}
\label{eulerprodsection}
In this section, we will define the `Hecke operators at infinity', denoted as $T_{p,\cl{F}}^\infty$ (the notations will become clear later). We will prove some formalism for their actions on $f_\lambda$, and in the end show that the shape of the Euler product of the $L$-series implies that $f_\lambda$ is an eigenfunction of $T_{p,\cl{F}}^\infty$ for certain suitable $p$, and $\cl{F}$.

Throughout this section, we will fix the following data unless stated otherwise:
\begin{enumerate}[label=(\arabic*)]
    \item a fixed $\lambda$ with $1\leq \lambda\leq h$;
    \item a prime ideal $\fr{p}$ with $[\fr{p}]=[\clo_F]$ in $\Cl^+(F)$ and $\fr{p}\nmid\fr{n}\fr{d}^2$;
    \item a totally positive generator $p$ of $\fr{p}$;
    \item a set $\cl{F}$ of representatives of $\clo_F/\fr{p}$.
\end{enumerate}
From now on, we will always write `$(\fr{p},p,\cl{F})$ as before' whenever we mean a triple $(\fr{p},p,\cl{F})$ satisfying (2), (3), and (4) as above.

For a holomorphic function $g:\clh^r\rightarrow\comp$, define
\begin{align*}
    T_{p,\cl{F}}^\infty\thin g\deq \N(\fr{p})^{k_0/2-1}\bigg(g|_\mbf{k}\m{p&0\\0&1}+\sum_{\alpha\in\cl{F}}g|_\mbf{k}\m{1&\alpha\\0&p}\bigg).
\end{align*}
Explicitly, we have
\begin{align}
\label{explicithecke}
    (T_{p,\cl{F}}^\infty \thin g)(z)\eq p^{\mbf{k}+(\mbf{k}^\prime/2)-\bone}g(pz)+p^{\mbf{k}^\prime/2-\bone}\sum_{\alpha\in\cl{F}}g\big(\tfrac{z+\alpha}{p}\big),
\end{align}
where $\bone:=(1,1,\cdots,1)\in\inte^r$. 

Recall that $f_\lambda:\clh^r\rightarrow\comp$ is given by
\begin{align*}
    f_\lambda(z)\eq\sum_{\xi\in(\fr{t}_\lambda\fr{d}^{-1})_+}a_\lambda(\xi)\thin e(\tr(\xi z)),
\end{align*}
where $a_\lambda(\xi)=\N(\fr{t}_\lambda)^{k_0/2}A(\xi\fr{t}_\lambda^{-1}\fr{d})\xi^{-\mbf{k}^\prime/2}$ for $\xi\in(\fr{t}_\lambda\fr{d}^{-1})_+$.

\begin{prop}
The Fourier series expansion of $T_{p,\cl{F}}^\infty f_\lambda$ is
\begin{align*}
    (T_{p,\cl{F}}^\infty\thin f_\lambda)(z)\eq\sum_{\xi\in(\fr{t}_\lambda\fr{d}^{-1})_+}b(\xi)\thin e(\tr(\xi z)),
\end{align*}
where
\begin{align*}
    b(\xi)\eq p^{\mbf{k}^\prime/2}\thin \big(a_\lambda(\xi p)+p^{\mbf{k}-\bone}\thin a_\lambda\big(\tfrac{\xi}{p}\big)\big)
\end{align*}
for $\xi\in(\fr{t}_\lambda\fr{d}^{-1})_+$, and $a_\lambda(\alpha):=0$ for $\alpha\notin(\fr{t}_\lambda\fr{d}^{-1})_+$ by convention.
\end{prop}
\begin{proof}
We have
\begin{align*}
    p^{\mbf{k}+(\mbf{k}^\prime/2)-\bone}f_\lambda(pz)\eq p^{\mbf{k}+(\mbf{k}^\prime/2)-\bone}\;\sum_{\mathclap{\xi\in(\fr{t}_\lambda\fr{d}^{-1})_+}}\;a_\lambda(\xi)\thin e(\tr(\xi pz))\eq p^{\mbf{k}+(\mbf{k}^\prime/2)-\bone}\;\sum_{\mathclap{\xi\in(\fr{t}_\lambda\fr{d}^{-1})_+}}\;a_\lambda\big(\tfrac{\xi}{p}\big)\thin e(\tr(\xi z))
\end{align*}
and
\begin{align*}
    p^{\mbf{k}^\prime/2-\bone}\sum_{\alpha\in\cl{F}}f\bigg(\frac{z+\alpha}{p}\bigg)&\eq p^{\mbf{k}^\prime/2-\bone}\sum_{\alpha\in\cl{F}}\sum_{\xi\in(\fr{t}_\lambda\fr{d}^{-1})_+}a_\lambda(\xi)\thin e\big(\tr\big(\tfrac{\xi(z+\alpha)}{p}\big)\big) \\
    &\eq p^{\mbf{k}^\prime/2-\bone}\sum_{\xi\in(\fr{t}_\lambda\fr{d}^{-1})_+}\bigg(\sum_{\alpha\in\cl{F}}e\big(\tr\big(\tfrac{\xi \alpha}{p}\big)\big)\bigg)\thin a_\lambda(\xi)\thin e\big(\tr\big(\tfrac{\xi z}{p}\big)\big) \\
    &\eq p^{\mbf{k}^\prime/2-\bone}\sum_{\substack{\xi\in(\fr{t}_\lambda\fr{d}^{-1})_+\\ \xi/p\in(\fr{t}_\lambda\fr{d}^{-1})_+}}|\clo_F/\fr{p}|\cdot a_\lambda(\xi)\thin e\big(\tr\big(\tfrac{\xi z}{p}\big)\big) \\
    &\eq p^{\mbf{k}^\prime/2} \sum_{\xi\in(\fr{t}_\lambda\fr{d}^{-1})_+}a_\lambda(\xi p)\thin e(\tr(\xi z))
\end{align*}
where in the third line we use that
\begin{align*}
    \sum_{\alpha\in\cl{F}}e\big(\tr\big(\tfrac{\xi \alpha}{p}\big)\big)\eq\begin{cases} |\clo_F/\fr{p}|&\text{ if }\xi/p\in(\fr{t}_\lambda\fr{d}^{-1})_+; \\ 0&\text{ if }\xi/p\notin(\fr{t}_\lambda\fr{d}^{-1})_+, 
    \end{cases}
\end{align*}
which follows from viewing $\alpha\mapsto e(\tr(\xi\alpha/p))$ as a character on $\clo_F/\fr{p}$ since $\xi\in\fr{d}^{-1}$. The result then follows from \Cref{explicithecke}.
\end{proof}

\begin{prop}
\label{eigenfourier}
The following are equivalent:
\begin{enumerate}[label=\rm{(\roman*)}]
    \item $T_{p,\cl{F}} f_\lambda=A(\fr{p})f_\lambda$;
    \item $A(\fr{a})A(\fr{p})=A(\fr{a}\fr{p})+\N(\fr{p})^{k_0-1}A(\fr{a}/\fr{p})$ for all nonzero integral ideals $\fr{a}\in\cl{C}_\lambda=[\fr{t}_\lambda^{-1}\fr{d}]$, where \mbox{$A(\fr{b})=0$} if $\fr{b}$ is not an integral ideal by convention.
\end{enumerate}
\end{prop}
\begin{proof}
Write $b(\xi)$ as the Fourier coefficients of $T_{p,\cl{F}}^\infty f_\lambda$ as before. Note that (i) is equivalent to that for all $\xi\in(\fr{t}_\lambda\fr{d}^{-1})_+$,
\begin{align*}
    A(\fr{p})\thin a_\lambda(\xi)\eq b(\xi)\eq p^{\mbf{k}^\prime/2}\thin \big(a_\lambda(\xi p)+p^{\mbf{k}-\bone}\thin a_\lambda\big(\tfrac{\xi}{p}\big)\big).
\end{align*}
Using $a_\lambda(\xi)=\N(\fr{t}_\lambda)^{k_0/2}A(\xi\fr{t}_\lambda^{-1}\fr{d})\xi^{-\mbf{k}^\prime/2}$, this is equivalent to that for all $\xi\in(\fr{t}_\lambda\fr{d}^{-1})_+$,
\begin{align*}
    A(\fr{p})A(\xi\fr{t}_\lambda^{-1}\fr{d})\eq A((\xi\fr{t}_\lambda^{-1}\fr{d})\fr{p})+\N(\fr{p})^{k_0-1}A\big((\xi\fr{t}_\lambda^{-1}\fr{d})/\fr{p}\big),
\end{align*}
which is then equivalent to (ii) by setting $\fr{a}=\xi\fr{t}_\lambda^{-1}\fr{d}$.
\end{proof}

Now, we will go back to the Euler product of the $L$-series $L(s)$. Suppose that $A(\clo_F)=1$ and that $L(s)$ has a $\fr{p}$-partial Euler product as the following
\begin{align*}
    L(s)\eq L^\fr{p}(s)\;L_\fr{p}(s)
\end{align*}
where
\begin{align*}
    L_\fr{p}(s)\eq\sum_{j=0}^{\infty}A(\fr{p}^j)\cdot\N(\fr{p}^j)^{-s}\hspace{0.5cm}\text{and}\hspace{0.5cm}
    L^\fr{p}(s)\eq\sum_{\fr{p}\nmid\fr{a}}A(\fr{a})\cdot\N(\fr{a})^{-s}.
\end{align*}
Let $\cl{C}\in\Cl^+(F)$ be an equivalence class. Since $[\fr{p}]=[\clo_F]$ in $\Cl^+(F)$, the partial $L$-series $L(s,\cl{C})$ also has a $\fr{p}$-partial Euler product
\begin{align*}
    L(s,\cl{C})\eq L^\fr{p}(s,\cl{C})\;L_\fr{p}(s,\cl{C})
\end{align*}
such that 
\begin{align*}
    L_\fr{p}(s,\cl{C})\eq L_\fr{p}(s)\hspace{0.5cm}\text{and}\hspace{0.5cm}
    L^\fr{p}(s,\cl{C})\eq\sum_{\substack{\fr{a}\in\cl{C}\\\fr{p}\nmid\fr{a}}}A(\fr{a})\cdot\N(\fr{a})^{-s}.
\end{align*}

\begin{prop}
\label{localfactoreigenfunction}
Fix $\lambda$ with $1\leq\lambda\leq h$. Let $\fr{p}$ be a prime ideal with $[\fr{p}]=[\clo_F]$ in $\Cl^+(F)$ and $\fr{p}\nmid\fr{n}\fr{d}^2$, $p$ be a totally positive generator of $\fr{p}$, and $\cl{F}$ be a set of representatives of $\clo_F/\fr{p}$. Then, the following are equivalent
\begin{enumerate}[label=\rm{(\roman*)}]
    \item $T_{p,\cl{F}}^\infty f_\lambda=A(\fr{p})f_\lambda$;
    \item $L_\fr{p}(s,\cl{C}_\lambda)=(1-A(\fr{p})\N(\fr{p})^{-s}+\N(\fr{p})^{k_0-1-2s})^{-1}$ where $\cl{C}_\lambda=[\fr{t}_\lambda^{-1}\fr{d}]$.
\end{enumerate}
\end{prop}
\begin{proof}
It is a standard computation that (ii) is equivalent to the condition (ii) in \Cref{eigenfourier}. The result then follows from \Cref{eigenfourier}.
\end{proof}

\begin{corol}
\label{eulercorol}
If the $L$-series $L(s)$ satisfies the Euler product as given in \Cref{mainthm}, then for each $1\leq\lambda\leq h$, $T_{p,\cl{F}}^\infty f_\lambda=A(\fr{p})f_\lambda$ for all $(\fr{p},p,\cl{F})$ as before.
\end{corol}

\section{Invariance of $f_\lambda$ under $\Gamma_0^+(\fr{t}_\lambda,\fr{n})$}
\label{gamma0-inv-sect}

In this section, we will prove the invariance of $f_\lambda$ under $\Gamma_0^+(\fr{t}_\lambda,\fr{n})$ under the assumptions on the functional equations and the Euler product given in \Cref{mainthm}. Throughout this section, we will always assume these two assumptions on the $L$-series $L(s)$. The proof mainly uses the idea from \cite[Theorem~2.2]{conjhecke}. The original idea dealt with elliptic modular forms and we modify it to apply it to Hilbert modular forms. As mentioned in the beginning, there is a small but easily fixable gap in the proof of \cite[Theorem~2.2]{conjhecke}, and we will provide our own fix at the end of this section. For the coherence of the discussion, we postpone one technical but important lemma to the end of this section, which is also closely related to our fix to the gap.

\begin{prop}
\label{gamma0-inv}
Let $\{\fr{t}_1,\cdots,\fr{t}_h\}$ be a set of representatives of $\Cl^+(F)$ such that each $\fr{t}_\lambda$ is a \textbf{prime} ideal. If the $L$-series $L(s)$ satisfies the assumptions on the analytic properties, the functional equations, and the Euler product in \Cref{mainthm}, then for each \mbox{$1\leq\lambda\leq h$}, the associated holomorphic function $f_\lambda$ satisfies that $f_\lambda|_\mbf{k}\gamma=f_\lambda$ for all $\gamma\in\Gamma_0^+(\fr{t}_\lambda,\fr{n})$.
\end{prop}

Before proceeding to the proof, we introduce one more notational convention which we will use in the proof. We will extend the slash operator to the group ring $\comp[\GL_2^+(F)]$ by setting
\begin{align*}
    g|_\mbf{k}\bigg(\sum_j a_j[\gamma_j]\bigg)\deq\sum_j a_j\thin g|_\mbf{k}\gamma_j
\end{align*}
for $a_j\in\comp$ and $\gamma_j\in\GL_2^+(F)$. Here we write $\gamma$ as an element in $\GL_2^+(F)$ and $[\gamma]$ as an element in $\comp[\GL_2^+(F)]$ to avoid confusion.
\\

Now, we go back to the proof. Recall that $\SG_0(\fr{t}_\lambda\fr{n})=\{\sm{a&b\\c&d}\in\SL_2(\clo_F)\mid c\in\fr{t}_\lambda\fr{n}\}$. We will first show that each $f_\lambda$ is invariant under $\SG_0(\fr{t}_\lambda\fr{n})$, for which we do not need to assume that each $\fr{t}_\lambda$ is a prime ideal.

Note that under the two assumptions in \Cref{mainthm}, we have obtained that for all $1\leq\lambda\leq h$,
\begin{enumerate}[label=\rm{(\alph*)}]
    \item $f_\lambda|_\mbf{k}T^\alpha=f_\lambda$ for all $\alpha\in\fr{t}_\lambda^{-1}$ (\Cref{basic});
    \item $f_\lambda|_\mbf{k}E_\eps=f_\lambda$ for all $\eps\in\clo_+^\times$ (\Cref{basic});
    \item $f_\lambda|_\mbf{k}W_{q_\lambda}=\epsilon f_{\wt{\lambda}}$ for any totally positive generator $q_\lambda$ of $\fr{t}_\lambda\fr{t}_{\wt{\lambda}}\fr{n}$ (\Cref{fqncorol});
    \item $f_\lambda|_\mbf{k}A_\beta=f_\lambda$ for all $\beta\in\fr{t}_\lambda\fr{n}$ (\Cref{fqncorol});
    \item $T_{p,\cl{F}}^\infty f_\lambda=A(\fr{p})f_\lambda$ for all $(\fr{p},p,\cl{F})$ as before (\Cref{eulercorol}).
\end{enumerate}
As $F$ is a totally real field of degree $r>1$, it follows from (a), (d), and \Cref{sgamma1-generator} that each $f_\lambda$ is already invariant under $\SG_1(\fr{t}_\lambda\fr{n})$, which is a finite index subgroup of $\SG_0(\fr{t}_\lambda\fr{n})$.

For the next lemma, we do a renormalization of the Hecke operators at infinity. Let $(\fr{p},p,\cl{F})$ be as before. Define
\begin{align*}
    \cl{T}_{p,\cl{F}}\deq\bm{p&0\\0&1}+\sum_{\alpha\in\cl{F}}\bm{1&\alpha\\0&p}.
\end{align*}
Indeed, $f_\lambda|_\mbf{k}\cl{T}_{p,\cl{F}}=\frac{A(\fr{p})}{\N(\fr{p})^{k_0/2-1}}f_\lambda$ by (e).

\begin{lemma}
\label{primesum}
Let $(\fr{p},p,\cl{F})$ be as before with $\fr{p}\nmid(q_\lambda)$. Let $\cl{F}^\times=\cl{F}\bslash\{0\}$. Then,
\begin{align*}
    \sum_{\alpha\in \cl{F}^\times}f_\lambda|_\mbf{k}([\gamma_{p,-\alpha}]-1)[T^{\alpha/p}]\eq 0,
\end{align*}
where $\gamma_{p,-\alpha}=\sm{p&-\alpha\\-\beta q_\lambda&\frac{\alpha\beta q_\lambda+1}{p}}$ and $\beta$ is the unique element in $\cl{F}^\times$ such that $p|\alpha\beta q_\lambda+1$.
\end{lemma}
\begin{proof}
Note that
\begin{align*}
    f_\lambda|_\mbf{k}[W_{q_\lambda}]\cl{T}_{p,\cl{F}}[W_{q_\lambda}^{-1}]\thin =\thin\epsilon f_{\wt{\lambda}}|_\mbf{k}\cl{T}_{p,\cl{F}}[W_{q_\lambda}^{-1}]\thin =\thin\frac{\epsilon A(\fr{p})}{\N(\fr{p})^{k_0/2-1}}f_{\wt{\lambda}}|_\mbf{k}[W_{q_\lambda}^{-1}]\thin =\thin\frac{A(\fr{p})}{\N(\fr{p})^{k_0/2-1}}f_\lambda\thin =\thin f_\lambda|_\mbf{k}\cl{T}_{p,\cl{F}}.
\end{align*}
Hence,
\begin{align*}
    0&\eq f_\lambda|_\mbf{k}[W_{q_\lambda}]\cl{T}_{p,\cl{F}}[W_{q_\lambda}^{-1}]-f_\lambda|_\mbf{k}\cl{T}_{p,\cl{F}} \\
    &\eq f_\lambda|_\mbf{k}\sum_{\beta\in F^\times}\bm{p&0\\-\beta q_\lambda&1}-f|_\mbf{k}\sum_{\alpha\in\cl{F}^\times}\bm{1&\alpha\\0&p} \\
    &\eq f_\lambda|_\mbf{k}\sum_{\alpha\in\cl{F}^\times}\bigg(\bm{p&-\alpha\\-\beta q_\lambda&\frac{\alpha\beta q_\lambda+1}{p}}-1\bigg)[T^{\alpha/p}]\bm{1&0\\0&p}.
\end{align*}
The result then follows.
\end{proof}

We will now begin to extend the group that $f_\lambda$ is invariant under from $\SG_1(\fr{t}_\lambda\fr{n})$ to $\SG_0(\fr{t}_\lambda\fr{n})$. This method does not reply on the exact shape of $\SG_1(\fr{t}_\lambda\fr{n})$ but rather that it is a finite index subgroup containing $T^\alpha$ and $A_\beta$ for all $\alpha\in\clo_F$ and $\beta\in\fr{t}_\lambda\fr{n}$. We will thus simply write $H=\SG_1(\fr{t}_\lambda\fr{n})$ throughout the proof. However, it may be possible to simplify the proof using ad hoc properties of $\SG_1(\fr{t}_\lambda\fr{n})$.

Let $u=[\SG_0(\fr{t}_\lambda\fr{n}):H]$ and choose $\{\gamma_1,\cdots,\gamma_u\}$ such that $\SG_0(\fr{t}_\lambda\fr{n})=\bigsqcup_{i=1}^u H\gamma_i$. We may replace $\gamma_i$ by~$A_\beta\gamma_i$ or $T^\alpha\gamma_i$ so that $\gamma_i=\sm{a_i&b_i\\c_i&d_i}$ with $a_i,c_i\neq0$. Let $H_i=\gamma_i^{-1}H\gamma_i\cap H$ and $g_i=f_\lambda|_\mbf{k}\gamma_i-f_\lambda$ for $1\leq i\leq u$. Then, $g_i$ is invariant under $H_i$ for $1\leq i\leq u$. Let $\{\delta_1,\cdots,\delta_r\}$ be a basis of $\clo_F$ over~$\inte$. For each $1\leq i\leq u$ and $1\leq l\leq r$, since \mbox{$[\SG_0(\fr{t}_\lambda\fr{n}):\gamma_i^{-1}H\gamma_i]=[\SG_0(\fr{t}_\lambda\fr{n}):H]<\infty$}, there exists $m_{il}\in\pos$ such that $\gamma_i^{-1}H\gamma_i(T^{\delta_l})^{m_{il}}=\gamma_i^{-1}H\gamma_i$. Set $m=\lcm_{1\leq i\leq u,1\leq l\leq r}\{m_{il}\}$. Then, \mbox{$\gamma_i^{-1}H\gamma_iT^{\alpha m}=\gamma_i^{-1}H\gamma_i$} for all $\alpha\in\clo_F$ and all $1\leq i\leq u$. This implies that $T^{\alpha m}\in H_i$ for all $1\leq i\leq u$. Hence, each $g_i$ has a Fourier series expansion
\begin{align*}
    \big(f_\lambda|_\mbf{k}([\gamma_i]-1)\big)(z)\eq g_i(z)\eq\sum_{\xi\in\fr{d}^{-1}}\mu_i(\xi)\thin e\big(\tr\big(\tfrac{\xi z}{m}\big)\big)
\end{align*}
for some $\mu_i(\xi)\in\comp$.

Let $(\fr{p},p,\cl{F})$ be as before with $\fr{p}\nmid(q_\lambda)$. Write $\cl{F}^\times=\cl{F}\bslash\{0\}$. Then, by \Cref{primesum},
\begin{align*}
    0&\eq\sum_{\alpha\in\cl{F}^\times}f_\lambda|_\mbf{k}([\gamma_{p,-\alpha}]-1)[T^{\alpha/p}] \\
    &\eq\sum_{i=1}^u\sum_{\substack{\alpha\in\cl{F}^\times \\ \gamma_{p,-\alpha}\in H\gamma_i}}f_\lambda|_\mbf{k}([\gamma_i]-1)[T^{\alpha/p}] \\
    &\eq\sum_{i=1}^u\sum_{\substack{\alpha\in\cl{F}^\times \\ \gamma_{p,-\alpha}\in H\gamma_i}}\sum_{\xi\in\fr{d}^{-1}}\mu_i(\xi)\; e\big(\tr\big(\tfrac{\xi(z+\alpha/p)}{m}\big)\big) \\
    &\eq\sum_{\xi\in\fr{d}^{-1}}\bigg(\sum_{i=1}^u\mu_i(\xi)\sum_{\substack{\alpha\in\cl{F}^\times \\ \gamma_{p,-\alpha}\in H\gamma_i}}e\big(\tr\big(\tfrac{\xi\alpha}{pm}\big)\big)\bigg)\thin e\big(\tr\big(\tfrac{\xi z}{m}\big)\big).
\end{align*}
This implies that for all $\xi\in\fr{d}^{-1}$,
\begin{align}
\label{lambdaexposum}
    \sum_{i=1}^u\mu_i(\xi)\bigg(\sum_{\substack{\alpha\in\cl{F}^\times \\ \gamma_{p,-\alpha}\in H\gamma_i}}e\big(\tr\big(\tfrac{\xi\alpha}{pm}\big)\big)\bigg)=0.
\end{align}

We will now prove that $\mu_i(\xi)=0$ for all $\xi\in\fr{d}^{-1}\bslash\{0\}$ so that $g_i$ is a constant for each $1\leq i\leq u$. Fix an element $\xi\in\fr{d}^{-1}\bslash\{0\}$. For each $1\leq j\leq u$, write $\gamma_j=\sm{a_j&b_j\\c_j&d_j}$ as before. Note that $a_j\neq 0$, $c_j\neq 0$, and $((a_j),(c_j))=1$ by our choice of $\gamma_j$. By \Cref{dirichlet}, we may choose $r_j\in\clo_F$ for each $1\leq j\leq u$ such that
\begin{enumerate}[label=(P\arabic*)]
    \item each $p_j=a_j+r_jc_j$ is totally positive and $(p_j)$ is a prime of degree $1$;
    \item $(p_j)\nmid mq_\lambda\N(\fr{d})\fr{n}$ and $(p_j)\nmid\xi\fr{d}$ for all $1\leq j\leq u$;
    \item all $\N(p_1),\cdots,\N(p_u)$ are distinct rational primes.
\end{enumerate}

Let $\alpha_j=-b_j-r_jd_j$. Note that $((p_j),(\alpha_j))=1$ and $(p_j)$ is a prime of degree~$1$, so by \Cref{standardfp}, $\cl{F}_j:=\{a\alpha_j\mid a\in\inte,1\leq a\leq\N(p_j)-1\}$ is a set of representatives of $(\clo_F/(p_j))^\times$. For $\alpha\in\cl{F}_j$, we write
\begin{align*}
    \gamma_{p_j,-\alpha}\eq\m{p_j&-\alpha\\-\beta q_\lambda&\frac{\alpha\beta q_\lambda+1}{p_j}}
\end{align*}
as before, where $\beta\in\cl{F}_j$ is the unique element such that $p_j|\alpha\beta q_\lambda+1$. Define
\begin{align*}
    \cl{F}_{i,j}\eq \{\alpha\in \cl{F}_j\mid \gamma_{p_j,-\alpha}\in H\gamma_i\}.
\end{align*}
Note that $T^{r_j}\gamma_j=\sm{p_j&-\alpha_j\\**&*}$ so there exists some $t_j\in\fr{t}_\lambda\fr{n}$ such that $\gamma_{p_j,-\alpha_j}=\sm{1&0\\t_j&1}T^{r_j}\gamma_j$.
Then, in particular, $\alpha_j\in \cl{F}_{j,j}$ as $\gamma_{p_j,-\alpha_j}=A_{t_j}T^{r_j}\gamma_j\in H\gamma_j$, so $\cl{F}_{j,j}\neq\emptyset$. Applying $p=p_j$ and $\cl{F}^\times=\cl{F}_j$ to \Cref{lambdaexposum}, we obtain that
\begin{align}
\label{lambdaexposumij}
    \sum_{i=1}^u\mu_i(\xi)\bigg(\sum_{\alpha\in \cl{F}_{i,j}}e\big(\tr\big(\tfrac{\xi\alpha}{p_jm}\big)\big)\bigg)\eq 0.
\end{align}
We will rewrite this equation in order to apply \Cref{nonzerodet}.

Let $q_j=\N(p_j)$ for each $1\leq j\leq u$. Define
\begin{align*}
    s_{i,j}\deq\{a\in\inte\mid 1\leq a\leq q_j-1, \;a\alpha_j\in \cl{F}_{i,j}\}\;\subseteq\;\{1,\cdots,q_j-1\}.
\end{align*}
Write $n_j=\tr\big(\xi\alpha_j\cdot\tfrac{q_j}{p_j}\big)\in\inte$. Note that $q_j\nmid m$ by (P2), and $q_j\nmid n_j$ by \Cref{npnmid} since $(p_j)\nmid \alpha_j\xi\fr{d}$ and $(p_j)\nmid(\N(\fr{d}))$. Then, for $a\in\{1,\cdots,q_j-1\}$, 
\begin{align*}
    \tr\bigg(\frac{\xi a\alpha_j}{p_jm}\bigg)\eq\frac{a}{mq_j}\cdot\tr\bigg(\xi\alpha_j \cdot\frac{q_j}{p_j}\bigg)\eq\frac{n_ja}{mq_j}
\end{align*}
and hence
\begin{align*}
    S_{i,j}\deq\sum_{\alpha\in \cl{F}_{i,j}}e\big(\tr\big(\tfrac{\xi\alpha}{p_jm}\big)\big)\eq\sum_{a\in s_{i,j}}e\big(\tr\big(\tfrac{\xi a\alpha_j}{p_jm}\big)\big)\eq\sum_{a\in s_{i,j}}e\big(\tfrac{n_ja}{mq_j}\big).
\end{align*}

To summarize, we obtain the data $(q_j,m_j=m,n_j,s_{i,j})_{i,j=1}^u$ which satisfies that
\begin{enumerate}[label=\rm{(\arabic*)}]
    \item all $q_j$ are distinct rational primes by (P1) and (P3);
    \item $q_j\nmid m_jn_j$ for all $j$ by (P1), (P2), and \Cref{npnmid};
    \item $q_j\nmid m_i$ for all $i$ and $j$ by (P2);
    \item $s_{i,j}\subseteq\{1,\cdots,q_j-1\}$ by the definition of $s_{i,j}$;
    \item $s_{i,j}\cap s_{i^\prime,j}=\emptyset$ for $i\neq i^\prime$ by the definition of $s_{i,j}$;
    \item $s_{i,i}\neq\emptyset$ for all $i$ as $\cl{F}_{i,i}\neq\emptyset$ for all $i$.
\end{enumerate}
Therefore, the data $(q_j,m_j=m,n_j,s_{i,j})_{i,j=1}^u$ satisfies the requirements in \Cref{nonzerodet}, so \mbox{$\det((S_{i,j})_{i,j=1}^u)\neq 0$}. By \Cref{lambdaexposumij},
\begin{align*}
    \sum_{i=1}^u\mu_i(\xi)\thin S_{i,j}\eq 0.
\end{align*}
for all $1\leq i,j\leq u$. This then implies that $\mu_i(\xi)=0$ for all $1\leq i\leq u$.

Note that $\xi\in\fr{d}^{-1}\bslash\{0\}$ is fixed but arbitrary. Hence, we obtain that $\mu_i(\xi)=0$ for all $\xi\in\fr{d}^{-1}\bslash\{0\}$ and all $1\leq i\leq u$. Thus, all $g_i$ are constant. Now, each $g_i$ is invariant under $H_i=\gamma_i^{-1}H\gamma_i\cap H$, which is a finite index subgroup of $\SL_2(\clo_F)$ since $H$ is. That is, each $g_i$ is a classical Hilbert modular form of weight $\mbf{k}$ on $H_i$. Since $\mbf{k}\in(2\inte^+)^r$ and $g_i$ is constant, $g_i$ must be zero. Hence, $g_i=0$ for all $1\leq i\leq u$, meaning that $f_\lambda|_\mbf{k}\gamma_i=f_\lambda$ for all $1\leq i\leq u$. Thus, $f_\lambda$ is invariant under $\SG_0(\fr{t}_\lambda\fr{n})$. 

To further extend the group to $\Gamma_0^+(\fr{t}_\lambda,\fr{n})$, we choose $\{\fr{t}_1,\cdots,\fr{t}_h\}$ such that each $\fr{t}_\lambda$ is a prime ideal. This is possible by considering the class field of $F$ which corresponds to $\Cl^+(F)$ and using Chebotarev's density theorem. Then, by \Cref{gamma0-generator}, $\Gamma_0^+(\fr{t}_\lambda,\fr{n})$ is generated by $\SG_0(\fr{t}_\lambda\fr{n})$, $T^\alpha$, and $E_\eps$ for all $\alpha\in\fr{t}_\lambda^{-1}$ and $\eps\in\clo_+^\times$. In particular, $f_\lambda$ is invariant under $\Gamma_0^+(\fr{t}_\lambda,\fr{n})$, so \Cref{gamma0-inv} follows.

In conclusion, under the assumptions on the functional equations and the Euler product in \Cref{mainthm}, $\mbf{f}=(f_1,\cdots,f_h)$ is an adelic Hilbert modular form of weight $\mbf{k}$ and level $K_0(\fr{n})$. It then follows that $L(s)=L(s,\mbf{f})$ is the $L$-function of a Hilbert modular form of weight $\mbf{k}$ and level $K_0(\fr{n})$. It is also easy to see that throughout the proof of \Cref{gamma0-inv}, we only use the local Euler factors at primes $\fr{p}$ that are of degree $1$, with $[\fr{p}]=[\clo_F]$ in $\Cl^+(F)$, and $\fr{p}\nmid\fr{n}\fr{d}^2$. Hence, the full Euler product in \Cref{mainthm} can be relaxed to a partial Euler product at those primes.
\\

Now, we introduce the technical lemma for showing that for fixed $\xi\in\fr{d}^{-1}\bslash\{0\}$, $\mu_i(\xi)=0$ for all $1\leq i\leq u$. This lemma is a slight generalization of \cite[Lemma~4.5]{conjhecke} and the proof also follows from the same idea.

\begin{lemma}
\label{nonzerodet}
Let $u\in\pos$. For each $1\leq j\leq u$, let $q_j$ be a rational prime, $n_j$ be an integer, and $m_j$ be a positive integer such that
\begin{enumerate}[label=\rm{(A\arabic*)}]
    \item all $q_j$ are distinct;
    \item $q_j\nmid m_jn_j$ for all $j$;
    \item $q_j\nmid m_i$ for all $i$ and $j$.
\end{enumerate}
Let $s_{i,j}\subseteq\{1,\cdots,q_j-1\}$ for each $1\leq i\leq u$ such that
\begin{enumerate}[label=\rm{(B\arabic*)}]
    \item $s_{i,j}\cap s_{i^\prime,j}=\emptyset$ for $i\neq i^\prime$;
    \item $s_{i,i}\neq\emptyset$ for all $i$.
\end{enumerate}
Write $S_{i,j}=\sum_{a\in s_{i,j}}e\big(\tfrac{n_ja}{m_jq_j}\big)$ for each $i,j$. Then, $\det((S_{i,j})_{i,j=1}^u)\neq 0$.
\end{lemma}
\begin{proof}
Note that replacing $m_j$ and $n_j$ with $\frac{m_j}{(m_j,n_j)}$ and $\frac{n_j}{(m_j,n_j)}$ does not affect the result, so we may assume that $(m_j,n_j)=1$ for all $1\leq j\leq u$. We also write $\zeta_j=e\big(\tfrac{n_j}{m_jq_j}\big)$ so that $S_{i,j}=\sum_{a\in s_{i,j}}\zeta_j^a$. Note that each $\zeta_j$ is a primitive $(m_jq_j)$-th root of unity by the coprime assumptions. We will now prove the result by induction on $u$.

Suppose that $u=1$. It suffices to show that $S_{1,1}\neq 0$. Write $Q(X)=\sum_{a\in s_{1,1}}X^{a-1}$. Then, $S_{1,1}=\zeta_1\cdot Q(\zeta_1)$, and $Q$~is a nonzero polynomial in $\rat[X]$ of degree $\leq q_1-2$ as \mbox{$\emptyset\neq s_{1,1}\subseteq\{1,\cdots,q_1-1\}$}. Note that
\begin{align*}
    [\rat(\zeta_1):\rat]\eq\vphi(m_1q_1)\eq\vphi(m_1)\vphi(q_1)\;\geq\;\vphi(q_1)\eq q_1-1\;>\;q_1-2.
\end{align*}
Since $\deg Q\leq q_1-2$ and $Q$ is nonzero, this implies that $Q(\zeta_1)\neq 0$, and hence $S_{1,1}\neq 0$.

Now, suppose that $u\geq 2$. Note that $S_{i,1}=\sum_{a\in s_{i,1}}\zeta_1^a$ so $\det((S_{i,j})_{i,j=1}^u)$ can be expressed as $P(\zeta_1)$ for some $P\in\rat(\zeta_2,\cdots,\zeta_u)[X]$ with $\deg P\leq q_1-1$. By (B1) and (B2), there exists an element $a$ such that $a\in s_{1,1}$ and $a\notin s_{i,1}$ for all $i\neq 1$, so the coefficient of the term $X^a$ in $P$ is the determinant of the matrix $(S_{i,j})_{i,j=2}^u$, which is nonzero by the induction hypothesis. Hence, $P$ is nonzero, and we may write $P(X)=X\cdot Q(X)$ for some nonzero $Q\in\rat(\zeta_2,\cdots,\zeta_u)[X]$ with $\deg Q\leq q_1-2$ as $s_{i,1}\subseteq\{1,\cdots,q_1-1\}$.

Write $m=\lcm(m_1,\cdots,m_u)$ and \mbox{$m^\prime=\lcm(m_2,\cdots,m_u)$}. By the coprime assumptions, 
\begin{align*}
    [\rat(\zeta_1,\cdots,\zeta_u):\rat(\zeta_2,\cdots,\zeta_u)]&\eq\frac{\vphi(mq_1\cdots q_u)}{\vphi(m^\prime q_2\cdots q_u)} \\
    &\eq \frac{\vphi(m)}{\vphi(m^\prime)}\cdot\vphi(q_1)\\
    &\;\geq\;\vphi(q_1)\eq q_1-1\;>\;q_1-2
\end{align*}
where $\vphi(m)/\vphi(m^\prime)\geq 1$ since $m^\prime|m$. Since $\deg Q\leq q_1-2$ and $Q$ is nonzero, this implies that $Q(\zeta_1)\neq 0$, and hence $\det((S_{i,j})_{i,j=1}^u)=P(\zeta_1)=\zeta_1\cdot Q(\zeta_1)\neq 0$.
\end{proof}

To finish this section, let us make a short discussion on the gap in the proof of \cite[Theorem~2.2]{conjhecke} and provide our own fix. The proof uses Dirichlet's theorem to obtain primes $q_i$ and integers $a_i$ such that $\gamma_{q_i,a_i}=\sm{q_i&a_i\\**&*}\in H g_i$ and then plugs these $\gamma_{q_i,a_i}$ into Equation (2.3) in their paper. However, these integers $a_i$ \textit{a priori} do not lie in $\{1,2,\cdots,q_i-1\}$ while Equation (2.3), as a corollary of Equation (2.1) in their paper, requires that $a_{il}\in\{1,2,\cdots,q_i-1\}$. It is thus not possible to simply plug in $\gamma_{q_i,a_i}$ into Equation (2.3).

The fix that we offer here is essentially what we do for the case of Hilbert modular forms. We first prove that Equation (2.1) still holds if we replace the summing set by any set of representatives $\cl{F}^\times$ of $(\inte/q\inte)^\times$. In their paper, this indeed follows from \cite[Lemma 4.1]{conjhecke}. Then, we choose primes $q_i$ and integers $a_i$ such that $\gamma_{q_i,a_i}\in H g_i$ as before, and for each $q_i$ we choose the set of representatives $\cl{F}_i=\{aa_i\mid 1\leq a\leq q_i-1\}$ of $(\inte/q_i\inte)^\times$. Now, we apply Equation (2.1) with $\cl{F}^\times=\cl{F}_i$ and we may proceed as what we have done in this section. That is, if we write
\begin{align*}
    s_{i,j}\eq\{a\in\inte\mid 1\leq a\leq q_j-1,\; \gamma_{q_j,aa_j}\in H g_j\},
\end{align*}
then the same argument yields a new version of Equation (2.4) in their proof as the following
\begin{align*}
    \sum_{i}\bigg(\sum_{a\in s_{i,j}}e\bigg(\frac{na_ja}{mq_j}\bigg)\bigg)\lambda_i(n)\eq 0 \hspace{0.3cm} \text{for all }j.
\end{align*}
To obtain a similar result on the nonvanishing of the determinant, we prove a modified version of \cite[Lemma~4.5]{conjhecke}, i.e., \Cref{nonzerodet} in this section. Now, we can apply \Cref{nonzerodet} in this section to the data $(q_j,m_j=m,n_j=na_j,s_{i,j})_{i,j}$ to deduce the nonvanishing of the determinant, and hence obtain that \mbox{$\lambda_i(n)=0$} for all $i$.


\appendix

\section{Lemmas on prime ideals}
\label{degree1}

In this appendix, we will prove all the lemmas we need on prime ideals. As these lemmas are general enough, we will simply be working with general number fields, and as before, we let $K$ denote a general number field to distinguish it from $F$. 

Before stating the first lemma, let us recall some definitions. For a set $S$ of primes in a number field $K$, define the \textit{natural density} of $S$ as
\begin{align*}
    \lim_{X\rightarrow\infty}\frac{\hash\{\fr{p}\text{ prime of }K\mid\N(\fr{p})\leq X,\;\fr{p}\in S\}}{\hash\{\fr{p}\text{ prime of }K\mid\N(\fr{p})\leq X\}}
\end{align*}
if it exists. For a prime $\fr{p}$ of $K$ and the rational prime $q$ lying under~$\fr{p}$, let $e(\fr{p})=e(\fr{p}/q)$ denote the ramification index and let $f(\fr{p})=f(\fr{p}/q)$ denote the residue degree. A prime $\fr{p}$ of $K$ is called \textit{of degree $1$} if it lies over a rational prime $q$ which splits completely in $K$, or equivalently, $e(\fr{p})\thin f(\fr{p})=1$.

\begin{lemma}
\label{deg1}
The set of primes of degree $1$ in a number field $K$ has density $1$.
\end{lemma}
\begin{proof}
This is a proof following \cite[Lemma~1]{ziegler}. Consider the set
\begin{align*}
    P\eq \{\fr{p}\text{ prime of }K\mid\N(\fr{p})\leq X,\; e(\fr{p})f(\fr{p})\geq2\}
\end{align*}
Note that $|\{\fr{p}\text{ of }K\mid\N(\fr{p})\leq X,\; e(\fr{p})\geq2\}|=O(1)$ as there are only finitely many primes of $\rat$ which ramify in $K$. Now, 
\begin{align*}
        |\{\fr{p}\text{ of }K\mid\N(\fr{p})\leq X,\; f(\fr{p})\geq2\}|&\;\leq\;[K:\rat]\cdot|\{q\text{ prime in }\rat\mid q^2\leq X\}| \\
        &\eq O(\pi(\sqrt{X}))  \\
        &\eq O(\sqrt{X}/\log X).
\end{align*}
Hence, the set $P$ has density $0$, so the result follows.
\end{proof}

\begin{lemma}
\label{local-quotient}
Let $K$ be a number field and $\fr{p}$ be a prime of $K$. Let $S\subseteq\clo_K-\fr{p}$ be a multiplicatively closed set. Then, the natural map $\clo_K/\fr{p}\rightarrow S^{-1}\clo_K/S^{-1}\fr{p}$ is an isomorphism.
\end{lemma}
\begin{proof}
Note that localization commutes with quotient (see \cite[\href{https://stacks.math.columbia.edu/tag/00CT}{Proposition 10.9.14 (00CT)}]{stacks-project}). That is, if we let $\ov{S}$ denote the image of $S$ under the quotient map $\clo_K\twoheadrightarrow\clo_K/\fr{p}$, then $\ov{S}^{-1}(\clo_K/\fr{p})\simarrow S^{-1}\clo_K/S^{-1}\fr{p}$. Note that $\ov{S}\subseteq(\clo_K/\fr{p})^\times$ since $S\subseteq\clo_K-\fr{p}$, so $\ov{S}^{-1}(\clo_K/\fr{p})=\clo_K/\fr{p}$.
\end{proof}

\begin{corol}
\label{ideal-local-quotient}
Let $K$ be a number field, $\fr{t}$ be a nonzero ideal of $K$, and $\fr{p}$ be a nonzero prime ideal of $K$ such that $\fr{p}\nmid(\N(\fr{t}))$. Then, the natural map $\clo_K/\fr{p}\rightarrow\fr{t}^{-1}/\fr{t}^{-1}\fr{p}$ is bijective.
\end{corol}
\begin{proof}
Write $T=\N(\fr{t})$ and let $S=\{T^n|n\in\nat\}$. As $\fr{p}\nmid (T)$, $S$ is a multiplicatively closed subset of $\clo_K-\fr{p}$, so by \Cref{local-quotient}, $\clo_K/\fr{p}\simarrow S^{-1}\clo_K/S^{-1}\fr{p}$. Also, note that $T\in\fr{t}$, so $T\xi\in\clo_K$ for all $\xi\in\fr{t}^{-1}$, and hence $\fr{t}^{-1}\subseteq S^{-1}\clo_K$. In particular, we obtain a sequence of maps
\begin{align*}
    \clo_K/\fr{p}\rightarrow \fr{t}^{-1}/\fr{t}^{-1}\fr{p}\rightarrow S^{-1}\clo_K/S^{-1}\fr{p}. 
\end{align*}
Since $\clo_K/\fr{p}\simarrow S^{-1}\clo_K/S^{-1}\fr{p}$, this forces the map $\clo_K/\fr{p}\rightarrow\fr{t}^{-1}/\fr{t}^{-1}\fr{p}$ to be bijective.
\end{proof}

Recall that an element $\alpha\in K$ is called \textit{totally positive} if $\sigma(\alpha)>0$ for all embeddings $\sigma:K\hookrightarrow\real$. In particular, if $\alpha\in K$ is totally positive, then $\N(\alpha)>0$. In the following three lemmas, $p\in\clo_K$ is a fixed totally positive element such that $(p)$ is a prime of degree $1$.

\begin{lemma}
\label{standardfp}
Let $\alpha\in\clo_K$ with $p\nmid\alpha$. Then, the set of representatives of~$\clo_K/(p)$ can be chosen as $\cl{F}_p^\alpha=\{a\alpha\mid a\in\inte, 0\leq a\leq \N(p)-1\}$.
\end{lemma}
\begin{proof}
Write $q=\N(p)$ a rational prime. For any $a,b\in \cl{F}_p^1$, $p|(a\alpha-b\alpha)$ implies $p|(a-b)$ and hence $\N(p)|\N(a-b)$, so $q|(a-b)^r$. As both $q$ and $a-b$ are integers and $q$ is a rational prime, this implies that $q|(a-b)$, and hence $a=b$ by the choice. It follows that $\cl{F}_p^\alpha$ is a set of representatives of $\clo_K/(p)$ as $|\cl{F}_p^\alpha|=\N(p)$.
\end{proof}

\begin{lemma}
\label{xi-represent}
Let $\fr{t}$ be a nonzero ideal of $\clo_K$ such that $p\nmid\N(\fr{t})$. Then, every $\xi\in\fr{t}^{-1}$ can be written as $\xi=n+p\eta$ for some $n\in\cl{F}_p^1=\{0,1,\cdots,\N(p)-1\}$ and $\eta\in\fr{t}^{-1}$.
\end{lemma}
\begin{proof}
Since $p\nmid\N(\fr{t})$, by \Cref{ideal-local-quotient}, the natural map $\clo_K/(p)\rightarrow\fr{t}^{-1}/\fr{t}^{-1}(p)$ is bijective. By \Cref{standardfp}, $\cl{F}_p^1$ is a set of representatives of $\clo_K/(p)$, so it is also a set of representatives of $\fr{t}^{-1}/\fr{t}^{-1}(p)$. That is, for any $\xi\in\fr{t}^{-1}$, there exist $n\in\cl{F}_p^1$ and $\eta\in\fr{t}^{-1}$ such that $\xi=n+p\eta$, which proves the result.
\end{proof}

\begin{lemma}
\label{npnmid}
Suppose that $p\nmid\N(\fr{d})$. Then, for any $\xi\in\fr{d}^{-1}$, $(p)|\thin\xi\fr{d}$ if and only if $\N(p)\thin|\thin\tr\big(\frac{\xi\N(p)}{p}\big)$.
\end{lemma}
\begin{proof}
First, $\tr\big(\frac{\xi\N(p)}{p}\big)\in\inte$ as $p\thin|\N(p)$. If $(p)|\thin\xi\fr{d}$, then $\xi/p\in\fr{d}^{-1}$, so $\tr\big(\frac{\xi\N(p)}{p}\big)=\N(p)\cdot\tr(\xi/p)$ is divisible by $\N(p)$. Now, suppose that $\N(p)\thin|\thin\tr\big(\frac{\xi\N(p)}{p}\big)$ so that $\tr\big(\frac{\xi}{p}\big)\in\inte$ and suppose also that $(p)\nmid\xi\fr{d}$.  By \Cref{xi-represent}, there exist $n\in\cl{F}_p^1$ and $\eta\in\fr{d}^{-1}$ such that $\xi=n+p\eta$. Then,
\begin{align*}
    \tr\bigg(\frac{\xi}{p}\bigg)\eq\tr\bigg(\frac{n+p\eta}{p}\bigg)\eq n\cdot\tr\bigg(\frac{1}{p}\bigg)+\tr(\eta).
\end{align*}
Here $\tr\big(\frac{\xi}{p}\big)\in\inte$ by the assumption and $\tr(\eta)\in\inte$ since $\eta\in\fr{d}^{-1}$. Hence, $n\cdot\tr(1/p)\in\inte$, where $n\neq 0$ since $(p)\nmid\xi\fr{d}$.

We will now prove that $\tr(\alpha/p)\in\inte$ for all $\alpha\in\clo_K$. First, for any $m\in\inte$, since $(n,\N(p))=1$, there exist $a,b\in\inte$ such that $m=an+b\thin\N(p)$. Hence,
\begin{align*}
    \tr\bigg(\frac{m}{p}\bigg)\eq\tr\bigg(\frac{an+b\thin\N(p)}{p}\bigg)\eq a\cdot n\cdot\tr\bigg(\frac{1}{p}\bigg)+b\cdot\tr\bigg(\frac{\N(p)}{p}\bigg)\in\inte.
\end{align*}
In particular, $\tr(m/p)\in\inte$ for all $m\in\cl{F}_p^1$. Now, as $\cl{F}_p^1$ is a set of representatives of $\clo_K/(p)$ by \Cref{standardfp}, each $\alpha\in\clo_K$ can be written as $\alpha=m+p\beta$ for some $m\in\cl{F}_p^1$ and $\beta\in\clo_K$. Hence,
\begin{align*}
    \tr\bigg(\frac{\alpha}{p}\bigg)\eq\tr\bigg(\frac{m+p\beta}{p}\bigg)\eq \tr\bigg(\frac{m}{p}\bigg)+\tr(\beta)\in\inte.
\end{align*}
Therefore, $\tr(\alpha/p)\in\inte$ for all $\alpha\in\clo_K$, so $1/p\in\fr{d}^{-1}$. In particular, $\frac{\N(\fr{d})}{p}\in\clo_K$ since $\N(\fr{d})\in\fr{d}$. This implies that $\N(p)|\N(\fr{d})^r$, and hence $p\thin|\N(\fr{d})^r$ since $p\thin|\N(p)$, which leads to contradiction as $p\nmid\N(\fr{d})$ by the assumption.
\end{proof}

The last two lemmas deal with Dirichlet's theorem on arithmetic progressions for general number fields. The next lemma is a fact on ray class groups that we will need for the proof of the last lemma.

\begin{lemma}
\label{rayclass}
Let $\fr{m}=\fr{m}_f\fr{m}_\infty$ be a modulus of $K$ where $\fr{m}_f$ is the non-archimedean part and $\fr{m}_\infty$ is the archimedean part. Let $C_\fr{m}$ be the ray class group for the modulus $\fr{m}$. Then, two nonzero integral ideals $\fr{a}$ and $\fr{b}$ correspond to the same class in $C_\fr{m}$ if and only if there exist nonzero $\alpha,\beta\in\clo_K$ such that
\begin{enumerate}[label=\rm{(\roman*)}]
    \item $\alpha\fr{a}=\beta\fr{b}$;
    \item $\alpha\equiv\beta\equiv1\pmod{\fr{m}_f}$;
    \item $\alpha$ and $\beta$ have the same sign for every real prime dividing $\fr{m}_\infty$.
\end{enumerate}
\end{lemma}
\begin{proof}
See Proposition 1.6 in \cite[Chapter~V]{milne}.
\end{proof}

\begin{lemma}
\label{dirichlet}
Let $a,c\in\clo_K$ be nonzero with $((a),(c))=1$. Then, there exist infinitely many totally positive elements $p\in\clo_K$ in the set \mbox{$\{a+nc\mid n\in\clo_K\}$} with distinct $(p)$ such that $(p)$ is a prime of degree $1$.
\end{lemma}
\begin{proof}
Note that $c|\N(c)$. Let $d\in\inte^+$ and $n=d\cdot\frac{|\N(c)|}{c}\in\clo_K$ so that $a+nc=a+d|\N(c)|$. By choosing $d$ large enough, we may replace $a$ by $a+d|\N(c)|$ to assume that $a$ is totally positive.

Let $\fr{m}=\fr{m}_f\fr{m}_\infty$ be a modulus of $K$ such that $\fr{m}_f=(c)$ and $\fr{m}_\infty$ is the product of all real places of $K$. Let $C_\fr{m}$ be the ray class group modulo~$\fr{m}$, let $L$ be the class field of $C_\fr{m}$, and let $\vphi:C_\fr{m}\rightarrow\Gal(L/K)$ be the Artin isomorphism. For a suitable ideal $\fr{a}$ of $\clo_K$, write $[\fr{a}]_\fr{m}$ as the equivalence class of $\fr{a}$ in $C_\fr{m}$. By Chebotarev's density theorem, the set of prime ideals $\fr{p}$ of $K$ satisfying that $\vphi([\fr{p}]_\fr{m})=\vphi([(a)]_\fr{m})$ (and hence $[\fr{p}]_\fr{m}=[(a)]_\fr{m}$) has a positive density among all prime ideals of $K$, where $[\fr{a}]_\fr{m}$ denotes the equivalence class of $\fr{a}$ in~$C_\fr{m}$. Since the set of prime ideals of degree $1$ in $K$ has density $1$ by \Cref{deg1}, this implies that there exist infinitely many prime ideals of $K$ of degree $1$ such that $[\fr{p}]_\fr{m}=[(a)]_\fr{m}$. 

Let $\fr{p}$ be a prime ideal of $K$ of degree $1$ with $[\fr{p}]_\fr{m}=[(a)]_\fr{m}$. By \Cref{rayclass}, there exist $\alpha,\beta\in\clo_K$ such that (i)~$\alpha(a)=\beta\fr{p}$, (ii) $\alpha\equiv\beta\equiv1\pmod{c}$, and (iii) $\alpha/\beta$ is totally positive. Since $\fr{p}$ and $(a)$ lie in the same equivalence class in $C_\fr{m}$, they also lie in the same equivalence class in $\Cl^+(K)$ under the surjection $C_\fr{m}\twoheadrightarrow C_{\fr{m}_\infty}=\Cl^+(K)$, and hence $\fr{p}$ is generated by some totally positive element $p\in\clo_K$. By~(i), there exist some $\eps\in\clo_K^\times$ such that $\alpha a=\eps\beta p$. Note that $\eps p=(\alpha/\beta)a$ is totally positive by (iii) and the assumption that $a$ is totally positive, so we may replace $p$ by $\eps p$ so that $\alpha a=\beta p$. By (ii), there exist $n_\alpha,n_\beta\in\clo_K$ such that $\alpha=1+n_\alpha c$ and $\beta=1+n_\beta c$. Rewriting $\alpha a=\beta p$ then implies that $a+(n_\alpha a-n_\beta p)c=p$. The result then follows.
\end{proof}

\bibliographystyle{amsalpha}
\bibliography{ref}

\end{document}
\newcommand{\etalchar}[1]{$^{#1}$}
\providecommand{\bysame}{\leavevmode\hbox to3em{\hrulefill}\thinspace}
\providecommand{\MR}{\relax\ifhmode\unskip\space\fi MR }
% \MRhref is called by the amsart/book/proc definition of \MR.
\providecommand{\MRhref}[2]{%
  \href{http://www.ams.org/mathscinet-getitem?mr=#1}{#2}
}
\providecommand{\href}[2]{#2}
\begin{thebibliography}{BBB{\etalchar{+}}18}

\bibitem[BBB{\etalchar{+}}18]{conjhecke}
Sandro Bettin, Jonathan~W. Bober, Andrew~R. Booker, Brian Conrey, Min Lee, Giuseppe Molteni, Thomas Oliver, David~J. Platt, and Raphael~S. Steiner, \emph{A conjectural extension of {H}ecke's converse theorem}, Ramanujan J. \textbf{47} (2018), no.~3, 659--684.

\bibitem[BK11]{bookerkrish-strengthen}
Andrew~R. Booker and M.~Krishnamurthy, \emph{A strengthening of the {${\rm GL}(2)$} converse theorem}, Compos. Math. \textbf{147} (2011), no.~3, 669--715.

\bibitem[BK13]{bookerkrish-refine}
\bysame, \emph{Further refinements of the {GL}(2) converse theorem}, Bull. Lond. Math. Soc. \textbf{45} (2013), no.~5, 987--1003.

\bibitem[BK16]{bookerkrish-gln}
\bysame, \emph{A converse theorem for {${\rm GL}(n)$}}, Adv. Math. \textbf{296} (2016), 154--180.

\bibitem[BMS67]{dedearith}
H.~Bass, J.~Milnor, and J.-P. Serre, \emph{Solution of the congruence subgroup problem for {${\rm SL}\sb{n}\,(n\geq 3)$} and {${\rm Sp}\sb{2n}\,(n\geq 2)$}}, Inst. Hautes \'{E}tudes Sci. Publ. Math. (1967), no.~33, 59--137.

\bibitem[Bru08]{123}
Jan~Hendrik Bruinier, \emph{Hilbert modular forms and their applications}, The 1-2-3 of modular forms, Universitext, Springer, Berlin, 2008, pp.~105--179.

\bibitem[Bum97]{bump}
Daniel Bump, \emph{Automorphic forms and representations}, Cambridge Studies in Advanced Mathematics, vol.~55, Cambridge University Press, Cambridge, 1997.

\bibitem[CF95]{conreyfarmer}
J.~B. Conrey and D.~W. Farmer, \emph{An extension of {H}ecke's converse theorem}, Internat. Math. Res. Notices (1995), no.~9, 445--463.

\bibitem[CPS94]{cogdellpsi}
J.~W. Cogdell and I.~I. Piatetski-Shapiro, \emph{Converse theorems for {${\rm GL}_n$}}, Inst. Hautes \'{E}tudes Sci. Publ. Math. (1994), no.~79, 157--214.

\bibitem[DN70]{doinaganuma}
Koji Doi and Hidehisa Naganuma, \emph{On the functional equation of certain {D}irichlet series}, Invent. Math. \textbf{9} (1969/70), 1--14.

\bibitem[Hec36]{hecke}
E.~Hecke, \emph{\"{U}ber die {B}estimmung {D}irichletscher {R}eihen durch ihre {F}unktionalgleichung}, Math. Ann. \textbf{112} (1936), no.~1, 664--699.

\bibitem[JL70]{jacquetlanglands}
H.~Jacquet and R.~P. Langlands, \emph{Automorphic forms on {${\rm GL}(2)$}}, Lecture Notes in Mathematics, vol. Vol. 114, Springer-Verlag, Berlin-New York, 1970.

\bibitem[Lie81]{liehl}
Bernhard Liehl, \emph{On the group {${\rm SL}\sb{2}$} over orders of arithmetic type}, J. Reine Angew. Math. \textbf{323} (1981), 153--171.

\bibitem[Mil20]{milne}
J.~S. Milne, \emph{Class field theory (v4.03)}, 2020, Available at \url{www.jmilne.org/math/}, pp.~287+viii.

\bibitem[Nag73]{naganuma}
Hidehisa Naganuma, \emph{On the coincidence of two {D}irichlet series associated with cusp forms of {H}ecke's ``{N}eben''-type and {H}ilbert modular forms over a real quadratic field}, J. Math. Soc. Japan \textbf{25} (1973), 547--555.

\bibitem[Neu99]{neukirch}
J\"{u}rgen Neukirch, \emph{Algebraic number theory}, Grundlehren der mathematischen Wissenschaften [Fundamental Principles of Mathematical Sciences], vol. 322, Springer-Verlag, Berlin, 1999, Translated from the 1992 German original and with a note by Norbert Schappacher, With a foreword by G. Harder.

\bibitem[PS75]{p-s}
I.~I. Piatetski-Shapiro, \emph{On the {W}eil-{J}acquet-{L}anglands theorem}, Lie groups and their representations ({P}roc. {S}ummer {S}chool, {B}olyai {J}\'{a}nos {M}ath. {S}oc., {B}udapest, 1971), Halsted Press, New York-Toronto, Ont., 1975, pp.~583--595.

\bibitem[Shi78]{shimurahmf}
Goro Shimura, \emph{The special values of the zeta functions associated with {H}ilbert modular forms}, Duke Math. J. \textbf{45} (1978), no.~3, 637--679.

\bibitem[{Sta}24]{stacks-project}
The {Stacks project authors}, \emph{The stacks project}, \url{https://stacks.math.columbia.edu}, 2024.

\bibitem[Vas72]{vaserstein}
L.~N. Vaserstein, \emph{The group {$SL\sb{2}$} over {D}edekind rings of arithmetic type}, Mat. Sb. (N.S.) \textbf{89(131)} (1972), 313--322, 351.

\bibitem[Wei67]{weil}
Andr\'{e} Weil, \emph{\"{U}ber die {B}estimmung {D}irichletscher {R}eihen durch {F}unktionalgleichungen}, Math. Ann. \textbf{168} (1967), 149--156.

\bibitem[Zie06]{ziegler}
Volker Ziegler, \emph{On the distribution of the order of number field elements modulo prime ideals}, Unif. Distrib. Theory \textbf{1} (2006), no.~1, 65--85.

\end{thebibliography}
